\chapter{\pysv\ 1.6--7}
\setlength{\footmarkwidth}{1.8em}
\pagewiselinenumbers


%[18,25]

\section{Sūtras 1.6 and 7 plus the Bhāṣya up to the definition of \emph{pratyakṣa}}
% This is found on Tm8r3 (right); L8r8 (L8r has 11 lines!)
\mll{BHINT6}
\bht{काः पुनस्ताः क्लिष्टाक्लिष्टाः पञ्चप्रकारा वृत्तय} इत्याह---\edn{don't forget the reference in BSBh 2.4.12: nanv atrāpi śrotrādinirapekṣā bhūtabhaviṣyadādiviṣayā aparā manaso vṛttir astīti samānaḥ pañcasaṃkhyātirekaḥ; evaṃ tarhi 'paramatam apratiṣiddham anumataṃ bhavati' iti nyāyāt ihāpi yogaśāstraprasiddhā manasaḥ pañca vṛttayaḥ parigṛhyante -- 'pramāṇaviparyayavikalpanidrāsmṛtayaḥ' nāma;}
\donote{BHINT6}{काः पुनस्ताः क्लिष्टाक्लिष्टाः पञ्चप्रकारा वृत्तयः}

% ya iti starts Tm8r4
\mll{STR6}%
\str{प्रमाणविपर्य्ययविकल्पनिद्रास्मृतय} इति॥६॥%
\donote{STR6}{\textbf{प्रमाणविपर्य्यय\- विकल्प\-निद्रास्मृतयः॥६॥}}

%[18,27]
\vrt{एतावत्य एव}{\Tm\Me}{\rdg{एतावत्येव}{L}} वृत्तयः। तत्र—\edn{should this tatra be part of the sūtra or the bhāṣya?}

\mll{STR7}%
\str{प्रत्यक्षानुमाना\vrtms{गमाः}{\Tmpc L\Me}{\rdg{॰गमा\il{}ः\ir{}}{\Tm}} प्रमाणानि॥७॥}%
\donote{STR7}{\textbf{प्रत्यक्षानुमानागमाः प्रमाणानि॥७॥}}%

\vrt{प्रमाणाख्या}{\eme}{\rdg{प्रमाख्या}{\Tm L\Me}} \mds वृत्तिस्त्रिधैव भिद्यते। तत्र प्रमा\mde णाख्यायाश्\vrt {चित्त\-वृ\-त्तेः}{\Tm\Lpc\Me}{\rdg{चित्तवृत्ते\il{}ः\ir{}}{L}} प्रथमो भेदः %(19,5)
प्रत्यक्षम्। अतस्तल्लक्षणमेवाभिधीयते तत्पूर्व्वक\-त्वाद्भेदद्वयस्य॥
\edn{That all the other pramāṇas presuppose pratyakṣa is not particularly special but it is interesting nonetheless that our author makes an explicit statement.}%

\mll{PRTYDEF}
\bht{इन्द्रिय\vrtms{प्रनाडिकय\mds ा}{L\Me}{%
	\rdg{प्रनाडिकाय्\lost{1}}{\Tmac},
	\rdg{प्रनाडिकय्\lost{1}}{\Tmpc}}}---%
इन्द्रियमिति \mde \vrt{श्रोत्रादि}{\Tm\Me}{%
	\rdg{श्रोत्राभि॰}{L}}
पञ्चकम्, न कर्म्मे\-न्द्रियम्, तस्य चित्तवृत्य\vrtmm{हेतुक}{\Tm L}{%
	\rdg{॰हेतु॰}{\Me}}%
\mds त्वात्।
\vrtsm{तत्}{L}{%
	\rdg{\om}{\Me}}%
क्रियार्त्थं हि तत्।
\edn{Perhaps tatkriyārthaṃ tat is a good reading. Cittavṛttyahetuka should refer to śrotrādi (they do not act because of cittavṛtti); "they (śrotrādi) exist for the sake of that (cittavṛtti)" remove "ka"! remove the first tat!}%
तस्मादिन्द्रिय\-शब्दस्य सा\-मान्य\mde ार्त्थत्वेऽपि \vrt{चित्तवृत्तेर्}{\Me}{%
	\rdg{उचितवृत्तेर्}{\Tm L}%,
	% \rdg{cittavṛtter}{\Me}%
	}\edn{tu and u are hardly distinguishable. This reading might be justified. However, it would be a little peculiar syntax...}
उप\mds व्याचिख्या\mde \vrtmm{सितत्वा}{\Tmpc L\Me}{%
	\rdg{॰सिसत्वा॰}{\Tmac}}%
च्छ्रोत्रादि\vrtms{विषयतैव}{\Tm\Me}{%
	\rdg{॰विषतयैव}{L}}\edn{The verb upavyākhyā appears in the ChU. Due probably to this, Śaṅkara uses this verb in a technical sense, it seems. Now I have a note on this peculiar verb.}%
। वृत्यर्त्थत्वाच्च श्रोत्रादी\mds नाम्।\edn{again, here appears to be a support that śrotrādi are there for the sake of cittavṛtti}
इन्द्रियमेव प्रनाडिका द्वारम्---%
शब्दाद्या\mde कार\-वृत्तिरूपेण परिणमम\mds ानस्य \bht{चित्तस्य}।
\vrt{अतस्}{L}{%
	\rdg{अतस्तत्}{\Me}}%
तेनेन्द्रियद्वारे\-ण \bht{सामान्य\-वि\-शेष\mde ात्मकबा\mds ह्यव\mde स्त्व}ा\mds कारतया परिणममा\mde नमुप\-र\-ज्यते।
तस्य तद्\-\bht{उप\-रागाद्}धेतोश्चित्तस्य या % L8v1 starts after this
मुद्रा\mds प्रतिमुद्रावद्\bht {वृत्तिः}। सामान्य\-वि\-शेषात्मक\-व\mde स्तूपरागेऽपि
\bht{%
	\vrtsm{विशेषाव\mds धारण}{L\Me}{%
		\rdg{वि\wmhl{1}एषाव\lost{4}}{\Tm}}%
	प्रधाना%
} सैव \bht{प्रत्यक्षम्प्रमाणाम्}॥% remove danda before?
\donote{PRTYDEF}{इन्द्रियप्रनाडिकया चित्तस्य सामान्यविशेषात्मकबाह्यवस्तूपरागाद्विशेषावधारणप्रधाना वृत्तिः प्रत्यक्षम्प्रमाणाम्॥}
%Bhāṣya: indriyapranāḍikayā cittasya sāmānyaviśeṣātmakabāhyavastūparāgād  viśeṣāvadhāraṇapradhānā vṛttiḥ pratyakṣam pramāṇām

\section{\emph{Pratyakṣa} and \emph{viśeṣa}}
\subsection{Why the Bhāṣya says \emph{pratyakṣa} is mainly concerned with difference}
%[19,13]
अथ स्वविषयसामान्यविशेषावभासकसाम\mde र्त्थ्ये श्रोत्रादीनाञ्चित्तस्य च
\vrt{किं\-कृतं}{L\Me}{%
	\rdg{किं\wmhl{1}तम्}{\Tm}}%
विशेषावधार\mds णप्राधान्यमिति---\edn{The choice of the verb avabhās and "ka"?}

उच्यते---विशेषप्रतिबन्धत्वाद्\mde\ धानो\-पादान\mds ादी\mde नां वि\mds
शेषावधारण\vrtms{प्र\-धा\-नेत्य्। आह---}{L}{%
	\rdg{प्रधानतेति (तेत्याह)}{\Me}}%
\edn{Something is lost, around here, I think. Something like "(viśeṣāvadhāraṇapradhāneti) yuktam| tadā kiṃ sāmānyan nāvadhāryata ity (āha)" Or simply replace āha with uktam? without the question mark; just divide by a comma? āha part of the previous sentence?}%
न सामान्यं नावधार्य्यते? गुणभूतन्तु गृह्यत एव। यथा नीलं रूपमिति
नील\mde ावधारण\vrtms{प्राधान्ये}{\Tmpc L\Me}{%
	\rdg{॰प्रधान्ये}{\Tmac}}
सत्यपि रूपमेव \vrt{नीलम्}{L\Me}{%
	\rdg{\wmhl{1}लम्}{\Tm}}
इति गुणभूतम\mds वधारयति तथा दुःखशब्दः सुखशब्दो मूढशब्द इ\mde ति च। \mds
% Tm8r left is over. This was the last line (line 9).
तथा च विशेषावधारणप्रधाना वृत्तिरिति ह्युपपद्यते॥

%[19,20]
किञ्च संशयविपर्य्ययनिमित्तत्वाच्च सामान्योपलब्धेः। % L8v4 starts after tvā
\edn{"kiñ ca" and "...tvāc ca" are redundant.}%
\vrtsm{सामान्ये ह्युप}{\Tmac\Tmpc L\Me}{%
	\rdg{सामान्ये \ccs ह्यु\cce \il{}ह्यु\ir{}प॰}{\Tm}}%
लब्धे बुभुत्सितविशेषस्य संशयो जा\mds यते।
% Tm8v1 after jā
समुपलब्धसामान्यस्मृतेर्व्वा विपर्य्ययः। समधिगतविशेषस्य तु \vrt{संशयविपर्ययौ}{\Me}{\rdg{संशयविपर्यया}{L}} न स्थितिमु%
% L8v5
पा\-श्नुवाते इति विशेषाव\mde धारणप्रध\mds ानतोच्यते॥
\edn{Sāmānya being the origin of saṃśaya and viparyaya could be a unique point of view. At least our author is following something. I should figure this out.}%

%[19,24]
यत्रापि सामान्यमेवावदिधारयिषितं गौरिति वाश्व \vrt{इति वा}{\Me}{%
	\rdg{इति वात्र}{L}}%
, तत्राप्यन्यत\-र\-परित्यागेनान्य\vrtms{तरावधारणेति}{\Me}{%
	\rdg{॰तराव\ilg धारणेति}{L}} % L8v6 starts after vi next line
विशेषावधारणप्रधानतैव। सर्वमे\-व हि \vrt{वस्तु}{\Lpc \Me}{%
	\rdg{वस्तु \ccs त\cce}{L}}
वस्त्वन्तरापेक्षया सामान्यं विशेषश्च॥

%[19,27] A new view
\subsection{Criticism of the view that \emph{sāmānya} is unreal (Buddhist)}
यस्य पुनः सामान्यमवस्त्वध्यारोपितमेवोषराम्बुवत्नेन्द्रियविषयमिति।
% L8v7 after ta next line
\edn{Cएर्तैन्ल्य Bउद्धिस्त्बुत् wहो? I wओउल्द्सय्रथेर् Dहर्मकीर्ति. Hए मदे थे पोइन्तर्थक्रियासमर्थं यत्तदत्र परमार्थसत् ।
अन्यत्संवृतिसत्प्रोक्तं ते स्वसामान्यलक्षणे ॥ प्रत्यक्ष 3 ॥}%
\edn{The matter is a bit more complicated. See the Tarkabhāṣā of Mokṣākaragupta toward the end of pratyakṣa and Kajiyama's notes 131ff (section 7.1; pp. 56ff.). According to Kajiyama, Dharmottara and Ratnakīrti contributed to Mokṣākara's saying indriya also takes sāmānya as object but I don't see those two clearly saying it in the passages cited by Kajiyama. The criticism that Buddhist won't acquire sāmānya to form inferences is already in Kumārila.}%
तस्य प्रतिनियत\-सामान्या\-ध्या\-रोपणा\-नुपपत्तिः। न ह्यनुपल\mde ब्धपू\-र्व्वमु\mds ख्ये \mde सम्भव\mds त्यध्यारोपणा, स्मृत्यसम्भवात्।
\edn{The meaning is clear but the use of the compound anupalabdhapūrvamukhya and its association of smṛti might be of interest}%
न कदाचिददृष्टमुख्योदकस्य सलिला\-ध्या\-रोपणमूषरेषु विद्यते॥
\edn{Is he talking about some sort of saltlake like places? The example used by the Buddhist means something like positing (salty) water in salt rock desert?}%

%[p]
अथास्मरन्नप्यध्यारोपयेत्॥
\edn{This, too, must have some background...}%

%[s]
रूपस्वलक्षणे शब्दत्वाद्यप्यध्यस्येत्॥
\edn{Again the meaning is more or less clear but apparently this is a very abridged form of a more lengthy discussion.}%

अथेन्द्रियान्तर\mde विषयत्वाद्\vrt {रूपस्वलक्ष\mds णे}{\Me}{%
	\rdg{रूपस्वलक्ष\lost{1}}{\Tm},
	\rdg{रूपस्वलक्षण॰}{L}}
\vrt{शब्दत्वाद्यध्यासो}{\Lpc\Me}{%
	\rdg{शब्दत्वा\ccs दु\cce द्यभ्यासो}{L}}
विरुद्ध इति चेत्---न, अवस्तुत्वे % L8v9 from saty
सत्यविशेषात्। \vrt{विशेषाभ्युपगमे}{\Lpc\Me}{\rdg{विशेषा\ccs भ्यु\cce भ्युपगमे}{L}} च वस्तुत्वप्रसंगः। श\-ब्दस्वलक्षणादिव रूपस्वलक्ष\mde णस्य॥
\edn{OK, here is more detail. Not necessarily an abliged discussion. Do we see something like this elsewhere?}%

%[20,6]
कि\mds ञ्चान्यत्---\vrtsm{देशकालान्यत्व}{\eme}{\rdg{देशकालान्य॰}{L\Me}}विशि\mde ष्ट\vrtms{प्रत्यक्षाभ्युपगमे}{\Tm L}{\rdg{प्रत्यक्षा[न]अभ्युपगमे}{\Me}}
\edn{There are only deśakālānyathā or deśakālānyatva in use it should be deśakālādi. See the next sentence. deśakālādyaviśiṣṭa...!}%
समस्तलोक\mds व्य\-व\-हार\vrtms{विलोप}{\Me({\sanskritfont ॰पः})}{%
	\rdg{॰विलोपा}{L}}
स्यात्।
%1 ) L8v10 after hī below
न हीदमहममुष्मिन्नवकाशे काले चास्मिन्नद्राक्षमेतस्मान्निमित्तादेवंविशेषणमिति च \vrt{स्मृत्युपपत्तिः।
% 2)
न}{\Me}{%
	\rdg{स्मृत्युपपत्तिन्न}{L}}
च \mde प्रत्यक्षा\vrtms{दृष्टे\-षु}{L\Me}{\rdg{॰दृ\wmhl{1}एषु}{\Tm}} विशेष्यवद्वि\vrtms{शेषणेषु}{L\Me}{\rdg{॰शेषणे\ccs नु\cce षु}{\Tm}} स्मृतिः संभवति।
% 3)
न च कल्प\mds नापोढेन व्य\-व\-हार\-\vrtms{गोचरातीतेन}{\Me}{%
	\rdg{गोचराती\wmhl{1}ए\wmhl{1}अ}{L}} % L8v11 at tī
व्यवहरमाणो दृश्यते।
\edn{conclusion p. 273 of Prasannapadā vol. 1.}%
तस्माल्लोकप्रसिद्धप्रत्यक्षानु\-सरण\-मेव न्याय्यम्॥

\section{\emph{Yogipratyakṣa} and \emph{mānasapratyakṣa}}
% [20,11]
\subsection{Question: Are they separate \emph{pramāṇa}(s) or kinds of \emph{pratyakṣa}?}
ननु च योगिनान्नि\mde र्व्विकल्पसमाधिजन्
\edn{This terminology reflects that of the Nyāyabhāṣya ./6\_sastra/3\_phil/nyaya/nyayabhasya.txt:ubhayathā darśanam, \il{}\il{}asty ātmā'' ity āptopadeśāt pratīyate, tatrānumānam --- \il{}\il{}icchādveṣaprayatnasukhaduḥkhajñānāny ātmano liṅgam'' iti, pratyakṣaṃ --- yuñjānasya yogasamādhijam \il{}\il{}ātmamanasoḥ saṃyogaviśeṣād ātmā pratyakṣa'' iti/ agnir āptopadeśāt pratīyate \il{}\il{}atrāgniḥ''iti, pratyāsīdatā dhūmadarśanenānumīyate, pratyāsannena ca pratyakṣata upalabhyate/ vyavasthā punaḥ --- \il{}\il{}agnihotra juhuyāt svargakāmaḥ'' iti, laukikasya svarge na liṅgadarśanam, na pratyakṣam/ stanayitnuśabde śrūyamāṇe śabdahetor anumānam, tatra na pratyakṣam, nāgamaḥ/ pāṇau pratyakṣata upalabhyamāne nānumānam, nāgama iti/}%
दर्शनन्तस्यानेवंलक्षणत्वात्तथा सुखरागादि%
\vrtmm{विज्ञान\mds स्य चानिन्द्रिय}{L\Me}{%
	\rdg{विज्ञान\lost{7}}{\Tm}}%
\vrtmm{प्रनाडिका}{L}{%
	\rdg{प्रनाडी॰}{\Me}}%
\mde पूर्व्वत्वात्\vrt {प्रमाणान्तरत्वम्}{\Tm\Me}{\rdg{प्रमाणान्तरम्}{L}} अ\-भ्य\-ुपेतव्यम्प्रत्यक्षत्वं वा वक्तव्यम्।

\subsection{Answer: \emph{Pratyakṣa}} % Change this at least!
% Consider not emending!
उच्यते---\vrt{न। प्रत्यक्षस्य}{\conj}{\rdg{न प्रमाणस्य}{\Tm L}, \rdg{प्रत्यक्षस्य (न प्रमाणस्य)}{\Me}} पुरुषप्रत्ययापेक्षत्वेन \vrt{परि\-हृत\-त्वात्।
न}{\Tmpc L\Me}{\rdg{परिहृतत्वेन}{\Tmac}} ह्यप्रत्यया \vrtsm{वृत्तिः प्र}{\Tm L\Me}{\rdg{वृत्तिः\ittspc{3}प्र॰}{\Tm}}त्यक्षम्।\edn{This part should mean "na pramāṇāntaram". Does this mean what may be called yogipratyakṣa or sukhaduḥkha, etc., are not pramāṇa because they do not [generate] a notion in people? And by puruṣapratyaya, does he mean something like hānopādāna? A tentative translation of this sentence would be "No. For, since pramāṇa is expected to [produce] notions [of hāna and upādān] in man, [they---yogins' darśana and sensations such as sukha-duḥkha] do not qualify [as pramāṇa]." Pratyaya in Śaṅkara is always in the form of a sentence? Not really. Even in the Vivaraṇa itself, there are usages like ahampratyaya or ihapratyaya.}\edn{20161014: The above was probably wrong. The answer "na" is probably OK. This should mean that there is no need to tell that those two are part of pratyakṣa or there is no need to postulate (an)other pramāṇa(s) to incorporate those. "parihṛtatvāt" means that those two possiblities are precluded because of the definition(?) that follows}\edn{Also, this part appears to tell a rathe unique view on pratyakṣa. Do the followings help? BSBh 1.3.33 pratyayāpratyayau hi sadbhāvāsadbhāvayoḥ kāraṇam; nānyārthatvam ananyārthatvaṃ vā; BSBh 2.2.25, 29?}
\edn{Does he mean anumāna, śabda/āptavacana, etc., to be apratyaya?}%
चैतन्योदय\vrtms{हेतोरेव}{\Tm\Lpc\Me}{%
	\rdg{हे\ccs तु\cce तोरेव}{L}}
सप्रत्य\-यायाः प्र\-त्य\mds क्षत्व\mde म्। % L9r2 starts at yāḥ pra
\edn{This sentence basically says the same thing as the previous one.}%
तथा चाह—%
\mll{PRMPHL}%
\bht{फलन्तदविशिष्टः पौरुषेयश्चित्त\-वृत्ति\-बोध} इति॥%
\donote{PRMPHL}{फलन्तदविशिष्टः पौरुषेयश्चित्त\-वृत्तिबोधः}

एवञ्च सति सुखरागादिज्ञानस्य
\edn{We could emend this to sukharāgādivijñānasya but may not be necessary. The two expressions mean the same thing.}%
\vrt{क्लिष्टाक्लिष्टरूपस्य}{\Tm\Me}{\rdg{क्लिष्टं क्लिष्टं रूपस्य}{L}}
\edn{kliṣṭākliṣṭa is an expression to qualify all the vṛttis. See YBh 1.3(?).}%
\vrtsm{तदविशिष्ट}{\Tmpc\Me L}{\rdg{तदविशिष्ट\ccs विशेष\cce}{\Tm}}पुरुषप्र\-त्ययफलावसानत्वात्प्रत्यक्षता \vrtsm{सिद्धै}{L\Me}{\rdg{सिद्धे॰}{\Tm}}व॥
\edn{Is the argument as follows?: The Bhāṣya says the result of pratyakṣa is tadaviśiṣṭaḥ pauruṣeyaś cittavṛttibodhaḥ; any cittavṛtti that produces that sort of result qualify as pratyakṣa; the cognition of sukha, rāga, etc., produce such results; therefore they are pratyakṣa? They are not another kind of pramāṇa because...they can be categorized in pratyakṣa?}%
% L9r3 ends with siddhaive


%[21,1]
इन्द्रियप्रनाडिकाग्रहणं लौकिक\vrtms{प्रत्यक्षोद्भावनानुवाद}{\Tm}{%
%	\rdg{pratyakṣotbhāvanānuvāda}{\Tm},
	\rdg{प्रत्यक्षोत्भावानुवाद}{L},
	\rdg{प्रत्यक्षोद्भवानुवाद}{\Me}}
एव। अन्यथा ही\-न्द्रिय\vrtms{निरपेक्षस्य}{\Tmpc L\Me}{%
	\rdg{निरपेक्ष\ccs त्व\cce स्य}{\Tm}}
\vrtsm{योगीश्वर}{L\Me}{%
	\rdg{\wmhl{2}गीश्वर॰}{\Tm}}%
\vrtmm{योर्ज्ञान}{\Tmpc L\Me}{%
	\rdg{योर्ज्जन॰}{\Tmac}}%
स्याप्रत्यक्षता स्यात्।
\edn{So, the first category in quesion, yogipratyakṣa, indeed is pratyakṣa.}%
% Tm9r starts at ādipratibandha
अस्ति हि क्लेशादि\vrtms{प्रतिबन्धाभावे}{\Tm\Lpc\Me}{%
	\rdg{॰प्रतिबन्धाभाव}{\Lac}}
सर्व्वार्त्थेन चित्तसत्वेन सर्व्वविषयेण % L9r4 after satvena sa
युगपत्सूक्ष्म\-व्यवहि\-ता\-दि\-ग्रहणम्। तथा चाह---\vrt{तारकं}{\Me}{\rdg{कारकम्}{\Tm L}}सर्वविषयं \vrtsm{सर्व्वथाविषयम}{\Tmpc\Me L}{\rdg{सर्व्वथाविषय\il{}म\ir{}॰}{\Tm}}क्रमञ्\vrt {चेति}{L\Me}{\rdg{चे\wmhl{1}}{\Tm}} \vrt{विवेकजञ्}{L\Me}{\rdg{वि\wmhl{1}एकजं}{\Tm}}ज्ञानमिति। \vrt{तस्मान्}{L\Me}{\rdg{\wmhl{1}स्मान्}{\Tm}}नावधार्य्यत इन्द्रियप्रनाडिकयै\-वेति॥
\edn{The use of iti at the end of this paragraph appears interesting.}%

%[21,6] L9r5 starts after bāhya
बाह्यवस्तूपरागादिति विपर्य्यय\vrtmm{निवृत्यर्त्थम्। विशेषा}{L\Me}{\rdg{॰निवृत्यर्त्थविशेषा॰}{\Tm}}वधारणप्रधान\-त्वेन संशयनिवृत्तिरिति॥
\edn{again iti at the end}%

\section{\emph{Pramāṇaphala}}
\subsection{Two views: \emph{pramāṇa} and its \emph{phala} are identical; there is an independent \emph{phala}}
%[21,8]
तत्र \vrt{केषाञ्चित्}{L\Me({\sanskritfont ॰आंचित्})}{\rdg{के\wmhl{2}ञ्चित्}{\Tm}}फलन्तदेव प्रमाणम्।\edn{I wonder if the author of the Bhāṣya in fact had Dignāga in mind. The order of the whole commentary (Vivaraṇa) on pratyakṣa reflects that of the Pramāṇasamuccaya. I wonder if this indeed reflects the thought of the Bhāṣyakāra. This seems unlikely. It was the Vivaraṇakāra's idea that somehow the definition in the Bhāṣya should cover yogipratyakṣa and mānasapratyakṣa. The Bhāṣyakāra apparently had no idea his definition would become problematic. Or did he?}
\edn{Is he basically following the Pramāṇasamuccaya? PS 1.8: savyāpārapratītatvāt pramāṇaṃ phalam eva sat.}%
अपरेषान्तु प्रमाणाद् \vrt{अन्यत्र}{\Tm L}{\rdg{अन्यत्(त्र)}{\Me}} \vrt{प्रमाणव्यक्त्यन्तरे}{\Tm L}{\rdg{प्रमाणव्यक्त्यन्तरम्(रे)}{\Me}}।
\edn{The above sentence looks really weird. I should be looking at the Tantravārttika, Prakaraṇapañcikā, NVTT, NBh, NMañ for hānopāānopekṣā. For pramāṇavyakti, see the TSP, NVTTP.}%
% L9r6 after gha below
\edn{Our author does not like either view.}%
यथा घटविज्ञानस्य प्रमाणास्य
\vrt{हानोपादानोपेक्षा}{\Me}{\rdg{हानोपादानापेक्षा}{\Tm}, \rdg{हानो\ccs चा\cce पाअदानोपेक्षा}{L}} बुद्धयः फलमिति
\edn{This is a quote from the Nyāyabhāṣya! No emendation???}%
\vrt{गुणदोषतदुभयशून्यत्वविषयत्वाद्}{\Me(\eme)}{%
  \rdg{गुणदोषवदुभयशून्यत्र\il{}त्व\ir{}विषय\-त्वाद्}{\Tm},
  \rdg{गुणदोषवदभवशून्यत्वविषयत्वाद्}{L},
  \rdg{गुणदोषतदुभय(वदभव)शून्यत्वविषय\-त्वाद्}{\Me}%
  }
उपादानादिबुद्धीनाम्%
\vrtsm{प्रमाणभूत}{\Tmpc L\Me}{\rdg{ब्रह्माणप्रत॰}{\Tmac}}घटवि\mds ज्ञानाद्
\vrt{अ\mde त्यन्तभेदः}{\Lpc\Me}{%
	\rdg{\lost{1}भ्यन्त\wmhl{1}ए\wmhl{2}}{\Tm},
	\rdg{अत्यन्त\ccs मेव\cce भेदः}{L}}%
॥
\edn{He is in fact following this Naiyāyika idea of hānopādānādibuddhis being the phala of pramāṇa? Most likely not. At the end, he says "nāpi tadvyatiriktapramāṇāntare."}%

\subsection{There is no independent \emph{pramāṇaphala}}
\vrtsm{अतश्च नान्य}{\eme}{\rdg{\wmhl{1}तश्चान्य॰}{\Tm}, \rdg{ततश्चान्य॰}{L\Me}}%
% L9r6 after vi below
\vrtms{विषयस्य}{\Tmpc L\Me}{\rdg{॰विषयीस्य}{\Tmac}}\edn{This should refer to another reason that is going to be stated why pramāṇa and pramāṇaphala are not two different things. Therefore, tataś ca (for the same reason) does not seem right. L9r6 after vi below}
\vrt{प्रमाणस्यान्यप्रमेयस्य}{}{%
  \rdg{प्रमाणस्यान्यप्रमेयस्य विषये}{\Tm L},
  \rdg{प्रमाणस्य अन्य(प्रमेयस्य)विषय॰}{\Me}
}
\edn{Could anyaprameyasya be an interpolation from a marginal comment?}%
\vrt{प्रमाणव्यक्त्यन्तरे}{\Tmpc}{%
  \rdg{प्रमाणव्यक्त्यन्तर॰}{\Tmac},
  \rdg{प्रमाणे व्यक्त्यन्तरे}{L},
  \rdg{॰प्रमाण(णे)व्यक्त्यन्तरं (रे)}{\Me}%
} फलमिति।
\edn{The whole of the above is hopelessly corrupt... but at least vyaktyantare phalam or pramāṇavykatyantare phalam appears consistent. So, that reading might be OK.}%
न हि खदिरच्छेदनस्य \vrt{खदिरोपादानन्}{\Lpc\Me}{%
	\rdg{खदिरोपादानान्}{\Tmac},
	\rdg{खदिरोपादा\ccs न\cce नन्}{L}}\vrt {तत्परित्यागो}{\Tm L\Me}{\rdg{तत्परित्यागो\ittspc{6}}{\Tm}} वा
\edn{The space left in Tm is curious. There could have been something written and erased although doing so without visible traces in a Malayalam manuscript should not be easy because writing literally leaves the surface scratched. A possible candidate that could come there is something like upekṣaṇa. There now is a research note on khadiracchedana. Nice parallels in the Upad and the BAUBh.}%
फलं \vrt{भवितुं}{\Me}{%
	\rdg{\ilg भवितु\il{}ं\ir{}}{\Tm},
	\rdg{भवतिम्}{\Lac}}
\edn{I cannot read what the letter before bha is in Tm; Td ignores it.}%
युक्तम्।
\vrtsm{वियोग}{\Me(\eme)}{%
	\rdg{अपि योग॰}{\Tmac},
	\rdg{अवियोग॰}{\Tmpc L},
	\rdg{वि(अपि)योग॰}{\Me}%
	}%
सम्बन्धफलत्वात्त्यागो\mds पादानयोः।
% L9r8 after tya
\edn{This sentence should explain why tyāga and upādāna are not the result of cutting of khadira tree but the result of ...yogasambandha. This might be simply viyogasambandhaphalatvāt, meaning that a relationship that is detached from khadiracchedana. The decision of whether taking or abondoning the resulting khadira wood comes not from the cutting but from somewhere else.}%
न च फलस्य फलं सम्भ\mde वति। \mds द्वै\mde धीभावफलव्यवहितत्वाच्च॥
\edn{NS/NBh/NV etc., BAUBh, ĀŚV, Bhāmatī, Upad etc. dvaidhībhāva means "making something into two." Something appears to have gone horribly wrong with these paragraphs. Cutting khadira is a Śābarabhāṣya thing. So, the translation might be something like "[Abondoning or taking of the khadira wood] is detached from the result that is making [the tree] into two."}%

% [21,16]
किञ्च फल\mds स्य \mde चेत्फलम्, क्रियानपेक्षाप्रसङ्गः। फलार्त्थं हि
\vrt{क्रिया\-पेक्ष्यते}{\Tmpc L\Me}{%
	\rdg{क्रियापेक्ष्यक्ष्यते}{\Tmac (\emph{kṣya} appearing twice at the end of a line and at the beginnning of the next line)}}%
। तच्चे\mds त्फ\mde लं विनापि क्रियया फलादपि
% L9r9 after kriyayā
\edn{phalād api part should mean something like "simply from the [first] result."}%
\vrtsm{सिध्यति, कोऽप्य}{\conj}{%
	\rdg{सिध्यति को ह्य॰}{\Tm},
	\rdg{सिद्ध्य को ह्य॰}{L}
	\rdg{(अतो ह्य) ततो न ह्य॰}{\Me}}%
नेकाय\mds ासा\-पेक्षितनिष्पत्तिक्रियामनुतिष्ठेत्।\edn{No need for conjecture???} \vrtsm{साधनान्वे\mde षणा}{\Me}{\rdg{साधनान्वेषणष{\cm\textbullet}आ॰}{L}}नुपपत्तिश्च, क्रि\mds याया अभा\-वा\mde त्॥

%[22,4]
नापि क्रियैव फलम्, \vrtsm{क्रियानुप}{L\Me}{\rdg{\wmhl{1}यानुप॰}{\Tm}}रम\vrtms{प्रसंगात्}{\Tmpc L\Me}{\rdg{॰प्रसंगा}{\Tmac}}।
\vrt{फलावसाना हि}{\eme}{\rdg{\ccs र्त्थपा\cce \il{}\ilg\ilg\ir{}वसानापि}{\Tm}, \rdg{फलावसानापि}{L}, \rdg{फलावसनापि [हि]}{\Me}} \mds क्रिया। क्रि\-यायाश्च
% L9r10 after phala below
फलत्वे फला\mde नभिलाषित्व\vrtms{प्रसङ्गः}{L\Me}{\rdg{॰प्रसं\wmhl{1}ः}{\Tm}}।
\mds तस्या दुःखत्वात्॥
\edn{Here it is again. The action being pain.}%

\subsection{\emph{Siddhānta}s}
%[22,6]
\subsubsection{No independent \emph{phala} nor are \emph{pramāṇa} and \emph{phala}  identical}
तस्मान्न प्रमाणस्वरूपमेव फलम्। \vrt{नापि}{\Me}{\rdg{नाधि}{L}} तद्व्य\mde ति\mds रिक्त\vrtms{प्रमाणान्तर}{L}{\rdg{प्रमाणान्तरम्}{\Me}} इत्याह---\bht{फल\mde न्तदविशिष्ट} इत्यादि। \vrt{तदविशिष्टस्}{L\Me}{\rdg{\wmhl{1}दविशिष्टस्}{\Tm}}तया \vrt{प्रमाणाख्यया}{\Tm\Lpc\Me}{\rdg{प्रमाणाख्याया}{\Lac}} वृ\mds त्या\-विशि\-ष्ट\-स्तुल्यस्तदाकारत्वेन।
% L9r10 after tadākā
\bht{पौ\mde रुषेयः} पुरुषस्यायम्\mds बोध इति। पुरुषस्य हि वृत्याकारता।
\edn{So far, it appears that our author is saying that in pratyakṣa (or in pramāṇa in general?), bodha and pramāṇa (a vṛtti) share the same ākāra, and the bodha and puruṣa share the same ākāra.}%

\subsubsection{The \emph{phala} is not the resulting \emph{dravya}, either}
तस्मान्\vrt {न क्रियाजन्}{\eme}{\rdg{नाक्रियाजं}{L\Me}}द्रव्यस्वरूपमात्रं फलं।

\subsubsection{The phala is a special state of the \emph{dravya}}

किन्तर्हि द्रव्यस्यै\mde
वा\vrtms{वस्थान्तरविशेषः क्रियाजः क्रियापर्यवसानकाले}{\eme}{%
	\rdg{॰वस्थान्तरवि\wmhl{2}ष\ccs क्रिया\cce\ins(कार्य्या)\ine जः \wmhl{1}र्य्यवसानकालः}{\Tm},
	\rdg{॰वस्थान्तरविशेषक्रियाजाक्रियापर्यावसानकालः}{L}, %
	\rdg{॰वस्थान्तरविशेषः (ष)क्रियाजः(जा)क्रियापर्यावसानकालः}{\Me}%
	}
क्रि\mds यो\vrtmm{परामोप}{L}{%
	\rdg{॰परमोप॰}{\Me}}%
पत्तेः फलम्।\edn{Is this a parallel? ./6\_sastra/3\_phil/vedanta/atmasiddhi.text:yastu sugatamatāvalambī vijñānābhinnahetujatayā tayorapi tadantarbhāvam abhimanyate; kaṇabhakṣapakṣāśrayaṇena vā tayorātmaviśeṣaguṇatvam;tābhyāṃ sukhaduḥkhādhikaraṇaṃ vyācakṣīta; svataḥsukhītyetadvimarśaṃ vātratyam / rāgadveṣādayastu caitanyasyaivāvasthāviśeṣāstadvadeva pratyakṣībhavantīti na tannidarśanenānumānodayaḥ / sukhaprayuktaviṣayīkāracaitanyaṃ rāgaḥ, tadvirodhaprayuktaviṣayīkāraṃ tadeva dveṣaḥ, bhūtaduḥkhajñānena cetaścalanaṃ śokaḥ, āgāmitajjñānena cetaścalanaṃ bhayamityādi lakṣaṇagranthādevāvagantavyamityalaṃ pravistareṇa / Or in BSBh, look for avasthāntara. It is the cause of bhedābhedabhāva Grammarians? About odana and phala} % L9v1 before tteḥ
फलनिष्पत्तावेव \mde हि \vrt{क्रियोपराम}{L}{%
	\rdg{षि\il{}(क्रि)\ir{}योपराम}{\Tm},
	\rdg{क्रियोपरम}{\Me}}
\mds उपपद्यते।
% bodhaḥ should be typeset in bold

\subsubsection{Either \emph{kriyā} or \emph{dravya} is figuratively the \emph{phala} of the other}
एवंभूतस्य मुख्यफलत्वे
\vrtms{प्रधानक्रियाव्यवहितविप्रकृष्टायाः}{\Lpc}{%
	\rdg{प्र\ccs मा\cce धानक्रियाव्यवहितविप्रकृष्टायाः}{L},
	\rdg{प्रधानक्रियाव्यवहित(विप्रकृष्टायाः)॰}{\Me}}
क्रियाया द्रव्यन्द्रव्यस्य वा क्रिया फलमित्य् \vrtsm{अ\mde न्योन्य}{L\Me}{%
	\rdg{\lost{1}न्यो\ins न्य\ine ॰}{\Tm}}% L9v2 after sthāpa
प्रत्युपस्था\vrtms{पकत्वेन}{L\Me}{%
	\rdg{॰पकत्व्\wmhl{1}न}{\Tm}}
प्र\-धान\-फलाभिमुखी\vrtms{करणात्}{\Tmpc L\Me}{%
	\rdg{करणा\ins त्\ine}{\Tm}}%
फ\mds लत्वोपचारो न विरोत्स्यते॥
\edn{So, basically the three things are not the phala: kriyā pramāṇa dravya. The last part I cannot say I fully understand but one can figuratively say either dravya or kriyā is the phala? I still don´t understand in what context our author is talking about the phala. Regular kriyā or specifically pramāṇa? Is there the picture of pramāṇa to kriyā to the change in dravya? Looks like pradhānakriyā is pramāṇa and pradhānakriyāvyavahitaviprakṛṣṭakriyā is the secondary;}%

\paragraph{Criticism: the \emph{phala} is only figurative?}
%[22,14] The following criticism appears reasonable. At any rate, what is not pramāṇaphala is more or less over.
ननु च भवतोऽप्युपचरितमेव फलम्, वृत्यविशिष्टायाः
\edn{buddheḥ? The bhāṣya actually reads "pauruṣeyo bodhaḥ." So, it should be masculine. Or he means vṛttyaviśiṣṭāyā avasthāyāḥ? <- That should be it!}%
\edn{Here is a nice one from the Upad: ./upad.txt:atrāha codakaḥ---avagatiḥ pramāṇānāṃ phalaṃ kūṭasthanityātmajyotiḥsvarūpeti ca vipratiṣiddham| ity uktavantam āha---na vipratiṣiddham| kathaṃ tarhi| kūṭasthanityāpi satī pratyakṣādipratyayānte lakṣyate, tādarthyāt| pratyakṣādipratyayasyānityatve 'nityeva bhavati| tena pramāṇānāṃ phalam ity upacaryate|108|}%
फलत्वेनाभ्युपगमात्। पुरुषस्य च \vrtsm{शुद्धत्वात्, तदाकारताया अनृत}{\Me}{%
	\rdg{शुद्धत्वादकारतायानृत॰}{L}}%
त्वम्। न हि स्फटिकमणेर% L9v3 after ma
\mde लक्तकोपधाना\vrtms{कारता}{L\Me}{\rdg{॰क्\wmhl{1}रता}{\Tm}} सत्या॥
\edn{upādhāna?}%

\paragraph{Answer: yes}
%[22,17]
बाढम्। तथापि त्वेतदेव %Tm9v1 starts here
फलम्मुख्यम्, दृष्टत्वात्, मुख्यफलान्तरादर्शनात्, सर्व्वस्यैव क्रियाकारकजातस्यैतत्फ\-लार्त्थत्वात्। एतदवसानं हि तत्सर्व्वम्।
% L9v4
तथा च प्रदर्शितम्॥

\paragraph{Criticism: then, the result being more or less erroneous, no reason to pursue \emph{pramāṇa}}
%[22,20]
ननु च \vrtsm{मिथ्या}{\Me}{\rdg{मित्थ्या॰}{\Tm L}}भूतस्यैव मुख्यफलत्वं वृत्याक\mde ारस्य,\edn{This notion that pramāṇa is anyway wrong is present in the intro to the BSBh.}
\edn{vṛtyākāra was puruṣa. This sentence does not make much sense.}%
तथा च तदभिलाषानुपपत्तिः।
\edn{A tentative translation: For [puruṣa] who takes the shape of the [citta]vṛtti, the main result [then] is nothing but unreal. That way, seeking [the main result of pratyakṣa/pramāṇa] will be impossible. doesn't tad refer to puruṣa?}%

\paragraph{Answer: correct}
बाढम्। \vrtsm{मिथ्या}{\Me}{\rdg{मित्थ्या॰}{\Tm L}}भूत % Tm9v2
\mds एव वृत्यनुगतो भोगः। तथा चाह---सत्वपुरुषयो\mde रत्यन्तासंकीर्ण\mds योः प्रत्ययाविशेषो भोग % L9v5 starts at rṇṇa
\edn{YS 3.35}%
 इति। भा\-ष्यमपि---तयोः स्वरूपोपलम्भे सति कुतो भोग इति॥
\edn{What he is doing in these paragraphs is that pramāṇa is essentially mithyā for mixing up puruṣa with cittavṛttis. Pratyakṣa/pramāṇa does not lead to mokṣa. It rather leads to saṃsāra.}%

\paragraph{\emph{samyagdarśana} (rather than \emph{pramāṇa}/\emph{pratyakṣa}) brings about \emph{nivṛtti}}
अत एव \vrt{च}{\Me}{\rdg{\om}{\Tm L}} सम्यग्दर्शनाद् \mde आत्यन्तिकी \vrt{निवृत्तिर्}{L\Me}{\rdg{नि\wmhl{1}त्तिर्}{\Tm}} उपपद्यते॥
\edn{If we adopt the reading na, samyagdarśana means pramāṇa, not vivekakhyāti. But the following suggest ca is better.}%

\paragraph{\emph{Mokṣa} is brought about only by \emph{samyagdarśana} (not by \emph{pramāṇa} that leads to action)}
%[22,25]
येषाम् \vrt{अमिथ्या}{L\Me}{\rdg{अमित्थ्या}{\Tm}} फलन्\vrt {ते\mds षां}{L\Me}{%
	\rdg{ते\lost{1}ए\lost{1}}{\Tmac}, % L9v6 before vairāgya
	\rdg{ते\lost{3}}{\Tmac}}
वैराग्यकारणाभावात्, सम्यग्दर्शना\mde न्वेषणानुपपत्ते\mds र्म्मोक्षाभावः। न च क्रियापरिनिष्पाद्यो मो\-क्षः, अनित्यत्वप्रसङ्गात्॥
\edn{This is such a typical Śaṅkara thing to say!}%

\subsection{The reason why \emph{puruṣa}’s nature, being aware of \emph{buddhi}, is not established here/Explanation of the intent of the last statement with regard to \emph{pratyakṣa} in the Bhāṣya}
%[23,4]
अथ बुद्धिबोद्धृत्वम्पुरुषस्य स\mde त्वं चासिद्धमिति चेत्,
% L9v7 after si above
आह---%
\mll{PRTSMVDBU}%
\bht{प्रतिसंवेदी बुद्धेः %Tm9v4
पु\mds रुष इत्युपरिष्टात्प्रवेदयिष्याम} इति॥\mde%
\donote{PRTSMVDBU}{प्रतिसंवेदी बुद्धेः पुरुष इत्युपरिष्टात्प्रवेदयिष्यामः}
\edn{This ends the discussion on pramāṇaphala. He does not believe in pramāṇa or pratyakṣa to lead to mokṣa.}%
\edn{The part being referred to in the Bhāṣya is probably not YBh 2.17 but YS 2.20. The Vivaraṇakāra refers to this very passage of YBh 1.7 when commenting on the sūtra (2.20) as the one being referred to be established in the future. YBh 2.20 also has "sa puruṣo buddheḥ pratisaṃvedī." In that sense, whether the reading was pravedayiṣyāmaḥ or upapādayiṣyāmaḥ may be determined by looking at MSS for the part where they record the Vivaraṇa on 2.20. The edition reads upapādayiṣyāmaḥ in YVi 2.20 (p. 189)... So, Td reads upapādayiṣyāma(ḥ) in YVi 2.20 (p. 161/PDF p. 61). Both readings are attested in different manuscripts. Does this mean that our author had access to more than one manuscript of the Yogabhāṣya?}%

\subsubsection{Criticism of the Naiyāyika position of \emph{pratyakṣa}, including their position on \emph{pramāṇaphala}}
% [23,6]
ये \edn{Naiyāyika definition of pratyakṣa; still inside the discussion on pramāṇaphala... Better consult NBh 1.1.4 or thereabout}चापीन्द्रियार्त्थसन्निक\mds र्षस्य प्रत्यक्षत्वमभ्युपयन्ति,
% L9v8 after kriya
न तेषामपि सम्यगभ्युपेतम्, सन्निकर्षस्य क्रियाफलत्वात्।
सन्निकर्ष\mde श्च संयोगः।
\vrt{स च}{L\Me}{\rdg{\wmhl{1}च}{\Tm}} क्रियाजस्तैरिष्यते।
यतो यदुपजा\mds यते \vrt{तत्तस्य}{\Me}{\rdg{तत्तस्य \ilg}{L}} प्रमाणम्। \vrt{त\mde च्च}{L\Me}{\rdg{त\ittspc{4}च्च}{L}} तस्य फलं यद् \vrt{उपजायते}{\Tm\Lpc\Me}{\rdg{उपजा\ccs\ilg\cce यते}{L}}॥
\edn{sannikarṣa = pratyakṣa sannikarṣa <- kriyā; sannikarṣa = kriyāphala sannikarṣa = saṃyoga saṃyoga = kriyāja (kriyāphala) i.e., sannikarṣa/saṃyoga comes from kriyā or kriyāphala, i.e., pratyakṣa is kriyāphala (the result of an action) A is born from B then B is A's pramāṇa; A that is born is the result (phala) of B B -> A then B is pramāṇa and A is phala This is a general principle that both sides should accept but according to the Naiyāyikas saṃyoga = sannikarṣa = pratyakṣa = pramāṇa is a phala of an action not the source (pramāṇa); here is a contradiction}%

% [23,9] L9v9 after iti
फलमपि \mds सन्निकर्षो ज्ञानम्प्रति प्रमाणमिति चेत्---
\edn{Then the Naiyāyikas say sannikarṣa = saṃyoga = pratyakṣa = pramāṇa is a phala (kriyāja) but toward jñāna, it is pramāṇa (source), generator}%

तच्च न,\edn{[Siddhāntin] The key here is the distinction between utpādaka and jñāpaka. Utpādaka should refer to sannikarṣa's being a kāraka} ज्ञ\mde ापकप्र\mds माण\mde प्र\mds स्तावे कारकप्रमा\mde णोपन्यासस्या\vrtms{युक्तत्वात्}{L\Me}{\rdg{॰यु\wmhl{1}\ccs स्य्\cce आत्}{\Tm}}। सन्निकर्षो हि ज्ञानस्योत्पाद\mds को \mde न तस्य ज्ञापकः।
% tasya should be understood as jñānasya; so, "na jñānasya jñāpakaḥ" L9v10 after kāraṇa
प्रमेयाधिगमाय हि प्र\-त्यक्षादीन्युप\mds न्यस्तानि, न ज्ञानोत्पत्तिकारणम्। \vrt{तत्प्रक्रिया\mde यां ह्यप्रकृतम्}{\conj}{%
	\rdg{\lost{4}यां ह्यप्रकृतं}{\Tm},
	\rdg{तत्प्रक्रियायान्यप्रकृतम्}{L},
	\rdg{तत्प्रक्रियायाम(यान्य)प्रकृतं}{\Me}} %
प्रक्रियेत॥
\edn{prakriyā is something like a themed discussion? Cf. prakaraṇa? The verb pra-kṛ means something like "to talk about something"?}%

%[23,13]
अथा\mds पि स्यात्\mde ---ज्ञानस्यापि प्रमेयत्वात्तच्चाप्युत्पद्यमानम्प्रमीयत इति तदुत्पत्तिकारणस्य सन्निकर्षस्य ज्ञापकत्वमेवेति॥
\edn{The point of this opponent's argument is jñānotpattikāraṇa is jñāpaka, and hence sannikarṣa is pramāṇa; no need to distinguish utpādaka and jñāpaka; makes sense, no?}%

तच्च न, ज्ञान\mds ं ज्ञनेनैव ज्ञायेत\mde, तेषान्तथाभ्युपगमात्।
% L9v11 after jñā of jñāyate above
\edn{Where do the Naiyāyikas teach this? The point should be that one cannot identify jñāpaka and jñānasyotpādaka because only another(?) jñāna is the jñāpaka of jñāna?}%
%
\edn{./6\_sastra/3\_phil/buddh/durvekamisra\_dharmottarapradipa.text:nanu kiṃ ghaḍādidṛṣṭāntabalāt prakāśāntaraprakāśyatā jñānasyāstām, āhosvit pradīpadṛṣṭāntasāmarthyāt prakāśāntaranirapekṣatayā svaprakāśatā? atrocyate | jñānasya svaprakāśarūpatvābhāve ghaṭādeḥ svarūpaprakāśanānupapatteḥ svaprakāśataiveṣṭavyā | atra ca prayogaḥ -- yad avyaktavyaktikaṃ na tad vyaktam yathā kiñcit kadācit kathañcid avyaktavyaktikam | avyaktavyaktikaś cāyaṃ ghaṭādiḥ jñānaparokṣatva iti vyāpakānupalabdhiprasaṅgaḥ | jñānasya jñānāntareṇa vyaktau hetur ayam asiddha iti cen na, nīlajñānodayakāle 'siddhatvād hetor nīlasya parokṣatvaprasañjanāt | na ca bhavatām api mate sarvaṃ vijñānam ekārthasamavāyinā jñānena jñāyate | bubhutsābhāve tadabhāvāt, yathopekṣaṇīyaviṣayā saṃvit | tat upekṣaṇīyam eva tāvad avyaktavyaktikatvād avyaktaṃ prasajyeta | na ceyam ./1\_veda/4\_upa/agamasastra\_w\_sankara.text:sopāyaṃ paramārthatattvaṃ laukikaṃ śuddhalaukikaṃ lokottaraṃ ca krameṇa yena jñānena jñāyate tajjñānam / ./1\_veda/4\_upa/bausbh.text:nanvanyajñānenānyanna jñāyata iti / ./6\_sastra/3\_phil/vedanta/bsbh.text:vyatireke hi satyekavijñānena sarvaṃ vijñāyata itīyaṃ pratijñā hrīyeta  /}%
न च सन्निकर्षो ज्ञानम्। \vrtsm{ज्ञानोत्पादक}{L\Me}{\rdg{ज्ञानो\wmhl{1}आदक॰}{\Tm}}त्वेनाभ्युपगमात्॥
\edn{So, jñāna, not being a jñāna, it cannot be known}%

% [23,16]
अथ \edn{This one sounds alright, again}ज्ञानस्य ज्ञाने सन्निकर्षो \vrt{निमित्तत्वात्प्रमाणम्}{\Me(conj.)}{\rdg{निमित्तप्रमाणत्वम्}{L\Tm}} इति \vrt{चेत्---%
अदृष्टे}{\Tmpc({\sanskritfont चेददृष्टे॰})\Me}{%
	\rdg{चे\il{}द\ir{}दृष्टे॰}{\Tm}, %
	\rdg{चेदष्टे}{L}}%
श्वरे\-च्छाविषयनिमेषोन्मेषादीनाम्\edn{Perhaps both jñānasya and jñāne are necessary in this context. No need to emend, probably} % L10r1 after ādīnām
अपि निमि\mds त्तत्वात्\mde प्रमाणत्वे कः प्रद्वेषः? न च ज्ञानस्य ज्ञेयत्व आत्मनश्च स्वभावचिद्रूपत्वानभ्युपगमे शक्यानवस्था परिहर्तुम्॥
\edn{"It is not possible to get rid of the anavasthā if jñāna is to be known and as long as one does not accept ātman being essentially the consciousness."}%

% [23,19]
अथात्मसमवेतमात्रतयैव % L10r2 after jñā next
ज्ञानस्य ज्ञेयता। न ज्ञानान्तरमपेक्ष्यते येनाति\vrtms{प्रसज्येतेति}{\Tm\Lpc\Me}{%
	\rdg{॰प्रसज्येते\ccs ते\cce ति}{L}} %
चेत्---

\vrt{तस्याप्य्}{\eme}{\rdg{अस्याप्य्}{\Tm L\Me}} अन्यत्प्रत्यक्षलक्षणं वाच्यम्। \vrtsm{अनिन्द्रियसन्निकर्षादिलक्षण}{L\Me}{\rdg{अनिन्द्रियस\wmhl{1}क\wmhl{5}ण}{\Tm}}त्वाद् \vrt{आत्मनि}{L\Me}{\rdg{आ\wmhl{2}}{\Tm}} ज्ञानसमवायत्वस्य। सन्निकर्षस्य च ज्ञानोत्पत्ति% This was the end of Tm9v L10r2 starts after jñāno
फलत्वात्तदु\-त्पत्तिव्यवधाने सति तस्य \vrt{ज्ञानस्य}{\eme}{\rdg{ज्ञानेन}{\Tm L\Me}} \vrt{नानन्तर्य्यम्}{\Tm\Me}{\rdg{\om}{L}}॥
% Was there the end of Tm9v somewhere in this paragraph?

% [23,23]
\vrtsm{तस्मात्तदा}{\Tm\Me}{%
	\rdg{\om}{L}}%
नन्\vrtms{तर्य्येणापि}{\Tm\Lpc\Me}{%
	\rdg{॰तर्य्यणापि}{\Lac}}
न सन्निक्र्षफलत्वं ज्ञानस्य। % Were these omissions supplied somewhere in L?
\edn{This ends the discussion on sannikarṣa}%

\paragraph{Examinations of the view that ideas of taking or removing are the result of \emph{pramāṇa}}
तथा हानादिबुद्धीनामपि \vrt{रूपादिज्ञानफलत्वानुपपत्तिरुक्ता}{L\Me}{\rdg{रूपादिज्ञाअ\wmhl{8}आ}{\Tm}}।
\edn{Well, the earlier discussion.}%
%[23,25] L10r4 after ca ya
व्यभिचाराच्च। \vrt{यथैकस्मिन्न्}{\Tm\Lpc\Me}{%
	\rdg{यथैवस्मिन्न्}{\Lac}}
\vrtsm{अंगनादि}{\Tm\Lpc\Me}{%
	\rdg{अंगानादि॰}{\Lac}}%
पिण्डविषये ज्ञाने समाने बहूनाम् \vrt{उत्पन्ने}{L\Me}{%
	\rdg{उल्प\ccs त्\cce न्न्\il{}ए\ir{}}{\Tm}}
कस्यचिदुपादानबुद्धिः, अन्यस्य हेयबुद्धिः, कस्यचिदुभयाभ\mds ावः\mde । तत्र हानोपादानबुद्धिर्व्व्यभिचरति।
\vrt{अन्यतरत्रैकमेव ज्ञानन्तद्}{L\Me}{% L10r5 before kam eva
	\rdg{अन्यतर\wmhl{1}ऐ\wmhl{6}आनन्तद्}{\Tm}}
\vrt{एवोभयं}{L\Me}{%
	\rdg{एवो\wmhl{1}यं}{\Tm}}
व्यभिचरत्युपेक्षकस्य। यच्च यद्व्यभिचरति न तत्तस्य फलम्॥
% t tasya phalaṃ starts Tm10r3.
\edn{There might be still textual problems but overall the idea is clear.}%

%[24,1]
किञ्च, हानोपादानाभावमात्रत्वादभाव एवोपेक्षा। न चाभावः \mds प्रमा\-ण\-स्य व\mde स्तुनः \mds % L10r6 after pha next line
फलम्\mde भवितुमर्हति। रागद्वेषाभ्याञ्चोपादानपरित्याग\-बु\-द्धी।
\vrt{तदभावात्}{L\Me}{%
	\rdg{तदभा\wmhl{1}त्}{\Tm}}%
\vrt{कारणाभावे}{\Tm\Lpc\Me}{%
	\rdg{कारणाभा\ccs बा\cce वे}{L}}
कार्य्याभाव इति प्रमाणविज्ञाने सत्य% Tm10r4 after sa
पि हा\-नो\-पादानबुद्ध्यभाव एवोपेक्षेति
% L10r6 after evope
व्यभिचारान्न हानादिबु\mds द्धिफलत्वम्॥
\edn{"...the notions of hāna, etc., are not phala" Again, the text looks more or less alright. There might be a significance with regard to our author's position on abhāva.... rāga -> upādāna(buddhi); dveṣa -> parityāga(buddhi)/hāna(buddhi)}%

%[24,5]
किञ्च, रागद्वेषतदभावा\mde नामप्युपादानादिबुद्धिफलत्वे प्रमाणत्व\-प्रसङ्गः।
\edn{Err, "There would be the undesirable consequence that rāga, dveṣa, and upekṣā are the pramāṇa with regard to the notions of upādāna, etc., being the result"? The idea seems OK given the diagram above.}%
पौरुषेयचित्तवृत्तिबोधफलत्वे तु न % Tm10r5 after ki L10r8 after kiñci
किञ्चिद्विरु\mds द्ध्यत इति॥
\edn{"But nothing contradicts if [the result/phala is] pauruṣeyacittavṛttibodha."??? iti is again significant here}%

\subsubsection{The problem of \emph{saṃvedyatva} of \emph{cittavṛtti}}
%[24,7] % XYZ
\paragraph{Is \emph{pramāṇa} a \emph{karaṇa}, one of \emph{kāraka}s?}

\mde ननु च चित्तवृत्तेः संवेद्यत्वमसिद्धम्। न हि क\mds रणमेव कर्म भवति। न हि दात्रमेव लूयते॥
\edn{Cittavṛtti and saṃvedyatvam; a new topic So, here cittavṛtti = pramāṇa = karaṇa; cittavṛtteḥ saṃvedyatvam is a problem because of the statement "phalaṃ tad(pramāṇākhyacittavṛtti)aviśiṣṭaḥ pauruṣeyaś cittavṛttibodhaḥ| pratisaṃvedī puruṣa ity upariṣṭāt pravedayiṣyāmaḥ" We have come back to the interpretation of the Bhāṣya For this opponent, the above implies that cittavṛtti is a karaṇa and it cannot be a karman}%

नै\mde ष दोषः। क्रियाकारक\mds योर्ब्भिन्नध\mde र्म्मत्वात्। क्रिया\vrtms{निर्वृत्तिम्}{\Me}{\rdg{निवृत्तिं}{\Tm}}
\edn{kriyānivṛtti or nirvṛtti? Rather nirvṛtti, no? The edition reads that so it is probably my typo when I made the e-text}%
प्रति स\-म्\-भू\-य % L10r9 after sambhūya
कारकाणि \vrt{व्याप्रियन्ते। % Tm10r6 after priya
न \mds परस्परं}{L}{%
	\rdg{व्याप्रियन्ते न \lost{5}}{\Tm},
	\rdg{व्याप्रियन्ते परस्परं}{\Me}}
विषयविषयिभावेन।  क्रि\mde या तु \vrt{सत्ता\-मा\-त्रेणैव}{\Me}{%
	\rdg{सत्तामा\wmhl{1}एण्\wmhl{2}व}{\Tm},
	\rdg{सत्तात्रेणैव}{L}}
\mds फलहेतुरिति करणत्वमस्यामुपचर्य्यते। न
\vrt{कारकवत्। कुर्व्वद्धि कारकम्। यदि क्रियापि कारकं स्या\mde न्न}{\conj}{%
	\rdg{[\mist{5}सर्व्व\mist{2}रक\mist{8}]\ccs स्व\cce आन्न}{\Tm},
	\rdg{कारकवत्कुर्व्वद्धि कारक यदि क्रियाधिकाक्रकं स्यान्न}{L},
	\rdg{कारकवत्। कुर्वद्धि कारयेदिति क्रियाधिकारकं स्यान्न}{\Me}}
भूतम्भव्यायोपदिश्येत॥ % L10r10 after bhavyā
\edn{This paragraph is totally obscure... Well, this is from the Śābarabhāṣya! Shows up in the ŚBh, NMañ, Bhāmatī, ŚV commentaries, etc. In essence, the author is saying that pramāṇa is kriyā and therefore not a kāraka and therefore not a karaṇa?}%

%[24,12]
\vrtsm{किञ्चानवस्था}{\Tmpc L\Me}{%
	\rdg{किञ्चनवस्था॰}{\Tmac}}%
प्रसंग\mds श्च। % Tm10r7
क्रिया चेत्करणं स्यात्, तन्निर्व्वर्त्त्यया \mde
\vrt{चान्य\-या क्रियया}{L\Me}{%
	\rdg{चा\wmhl{5}आ}{\Tm}} %
भवितव्य\mds म्। न ह्यन्यां क्रियामन्तरेण तस्याः
\vrt{करणत्वं}{\eme}{%
	\rdg{कारणत्वं}{L\Me}} %
घटते। यापि चान्या सापि
\vrt{करणं}{\eme}{%
%	\rdg{karaṇatvaṃ}{},
	\rdg{कारणं}{L\Me}} %
स्यादन्यस्याः क्रिय\mde ायाः। साप्यन्यस्या इत्यनवस्था। % L10r11 after s of sāpy
ततश्च क्रियासन्तानानुपर\mds मात्% Tm10r8
फलानभिनिर्व्वृत्तिः।
\vrt{तत्रान्या}{L}{%
	\rdg{तत्रान्त्या}{\Me}}
चेत्क्रियान्त\mde रेण विना फलहे\mds तुर्भवेत्, प्रथमस्य तथात्वे कः प्रद्वेषः।
\edn{This sentence does not look good, The last sentence should contribute to the point that kriyā is not a kāraka; and it appears that it is still in the context of anavasthā the reading antya, in contrast to prathama, might be good. These two paragraphs are really hard to figure out}%

\paragraph{\emph{Siddhānta}: \emph{cittavṛtti} is \emph{saṃvedya} as the revealer (\emph{abhivyañjaka})}

चित्तवृत्तेस्त्वभिव्यञ्जकत्वात्सम्वेद्यतेत्यदोषः।\edn{Does this part have anything to do with Nyāyasūtra Nyāyabhāṣya NS NBh 1.1.22?} % L10v1 starts after ado
\edn{"tu" is here. As opposed to what? Probably kriyā}%
%[24,18]
दृश्यते चा\mde भिव्यञ्जका\-नां
\vrtsm{स्वरूप}{L\Me}{%
	\rdg{स्वरू\wmhl{1}॰}{\Tm}}%
संवेद्यतयैव करणत्वम्प्र\mds दीपादीनाम्।
% Tm10r9 after dīpādīnām
\edn{pradīpa, etc., are abhivyañjaka; they are saṃvedya; by being naturally known, they are karaṇa; so is the cittavṛtti...}%
तथा चित्तवृत्तेरपि।

\subparagraph{Why is the eyesight not grasped while it is an \emph{abhivyañjaka}?}

चक्षुरादिष्वभि\mde व्यञ्जकत्वे \mds ऽप्यग्रहणाद्व्यभिचार
\edn{The sight, etc., do reveal things but they are not grasped, i.e., asaṃvedya, therefore there is a vyabhicāra: not everything that reveals does so by being saṃvedya.}%
इति \vrtsm{चेन्न, समान}{\conj}{%
	\rdg{चेतनसमान॰}{L}
	\rdg{चेत्, समान (चेतनसमान)}{\Me}}%
जातीयसंकरात्॥

कथम्?
\edn{The following should explain what he means by samānajātīyasaṃkarāt; why the sight, etc., are not saṃvedya "due to the confusion with things of the same kind"(?)}%

% \subparagraph{According to the theory that sense faculties are material}
%[24,21] L10v2 after yeṣa
येषान्तावद्भौतिकानीन्द्रियाणि,\edn{position A (Naiyāyikas)} तेषा\mde ञ्चक्षुषस्तैजसत्वाद्बाह्येनालोकेन \vrtsm{तुल्य}{\Tmpc L\Me}{\rdg{\il{}तु\ir{}ल्य॰}{\Tm}}जातीयेन
% Tm10v1 after jā
व्यतिभेदाद्गृह्यमाणस्यैवाविवेकः। यथा प्रदीपानां बहूनां प्रभाव्यतिभेदात्कस्य का प्रभेति % L10v3 afte prabheti
न विविच्यते। तथा श्रोत्रादिष्वपि॥
\edn{This paragraph appears to be relatively well-preserved}%

%[24,24]
% \subparagraph{According to the theory that sense faculties are not material}
येषामप्यभौतिकानि,\edn{Who?}
\vrt{तेषाम्}{\Me}{%
	\rdg{केषाम्}{L}}
अप्य\mde भिव्यञ्जकत्वेन समानजातीयता॥
\edn{Thus far is an explanation why cakṣur etc., are not saṃvedya even though they are abhivyañjaka}%

वृत्तेर्
\vrtsm{अप्य\mds व्यञ्जक}{\conj}{%
	\rdg{अप्यभिव्यञ्जक॰}{L},
	\rdg{अभिव्यञ्जक॰}{\Me}}%
त्वादग्रहणम्
\vrt{इति चेत्, तन्न}{L}{%
	\rdg{इति। तन्न}{\Me}}%
॥ % surviving part of Tm10v2 is below
\vrtsm{प्राप्तचैत\mde न्योपग्रह}{L\Me}{%
	\rdg{\lost{7}पग्रहण॰}{\Tm}}%
तय\mds ा
% L10v4 starts after vi of vibhāgya
\edn{The expression prāptacaitanyopagraharūpa is from YBh 2.20, 4.22. This adjective modifies buddhivṛtti The point the S has to prove is that a) cittavṛtti is abhivyañjaka and that b) it is not a grāhya...? It is not grāhya because it is a kriyā...? Shouldn't he be arguing that cittavṛtti is saṃvedya and it is a abhivyañjaka? The opponet here is saying "cittavṛtti is not a vyañjaka and that's why it is not cognized." Our author should counter on both counts, no?}%
\vrt{व्यञ्जकत्वाद्}{\conj}{%
	\rdg{विभाग्यत्वाद्}{L},
	\rdg{ऽभिभाव्यत्वाद्}{\Me}}
अव्यवहितत्वाच्च। सर्व्वगतत्वेऽप्यात्मनः सूक्ष्मत्वेन स्वच्छतरत्वेन च
\vrtsm{चित्तेना}{\Me}{%
	\rdg{चित्ते ता}{L}}%
व्यवधानाच्च ग्रहणं चित्त\mde वृत्तेः॥
\edn{Several things to note here. Puruṣa is equated as ātman, it is considered to be sarvagata and sūkṣma. This sentence should be explaninng how the grahaṇa (saṃvid) of cittavṛtti occurs and it appears as though it is not happening.}%

%%%% Conclusion!
\subsection{Conclusion: \emph{pramāṇaphala} is what the Bhāṣya says it is}
तस्माच्चित्तवृत्ति\vrtms{प्रमाणेनैव}{\Tmpc L\Me}{\rdg{॰प्रमाणैनैव}{\Tmac}} तदाकारतय\mds ा पौरुषेयोऽपि बोधः % L10v5 after bodhaḥ; Tm10v3 after taya
संवेद्य\-मान एव फलमिति सिद्धम् ॥
\edn{This must be the conclusion of the discussion on pramāṇaphala!}%

\subsubsection{How is \emph{cittavṛtti} object of awareness (\emph{saṃvedya})? Not as \emph{pratyakṣa}}
%[25,4]
येषां पुनः
\vrt{प्रत्यक्षेण संवेद्या}{\eme}{%
	\rdg{प्रत्यक्षेणासंवेद्या}{\all}}
चित्तवृत्तिस्%
\vrtsm{तेषाम्फल}{\Me}{%
	\rdg{तेषामफल}{L}}%
ग्रहणप्रमाणा\-भावः।
\edn{Who are these people? What's the point of bringing this up? This discussion in fact is related to the one that starts at XYZ Now I'm pretty sure that the reading should be pratyakṣeṇa saṃvedyā. The difference between the view being attacked and our author's is whether cittavṛtti is saṃvedya by being abhivyañjaka or by being pratyakṣa! Since it is hard to make sense out of this discussion, should we emend the text to "pratyakṣeṇa saṃvedyā"? That would be the Buddhist and Naiyāyika position. The last compound should mean the abhāva of both phalagrahaṇa and pramāṇa Should the last part be pramāṇavyāpārābhāvaḥ? So the whole compound is phalagrahaṇapramāṇavyāpārābhāvaḥ?}%
% This is Tm10v and L10v
फले च गृहीते \vrtsm{प्रमाणास्तित्व}{\eme}{%
	\rdg{ज्ञानास्तित्व}{L},
	\rdg{ज्ञानासिद्धत्व}{\Me}}%
\mde स्यार्त्थापत्तिः।
\edn{Is he talking about paracittavṛtti? Most unlikely. Arthāpatti suddenly appearing here is also suspicious. Maybe the argument is if phala is grasped, the existence of pramāṇa/jñāna is inevitable. But...}%
न च फलं गृह्य\-ते % L10v6
\vrtsm{प्रमाणव्या\-पारान}{\Tmpc L\Me}{%
	\rdg{प्रमाणव्यापाराद}{\Tmac}}%
भ्युपगमात्।
\edn{So, "ca" here means "but," stating what happens if we accept the opponent's position. The phala is not grasped and therefore the knowledge or pramāṇa does not exist. Because the opponent cannot accept the operation of pramāṇa because doing so would result in anavasthā situation.}%
%%% THE FOLLOWING IS IRRELEVANT This statement is not incompatible with the idea "pratyakṣeṇa saṃvedyā" because with that view, the final function of pramāṇa is not specified. cittavṛtti is still known by a pramāṇa but what happens then? A translation could be something like "Also, the effect is not grasped because [according to the theory] they do not accept the function of pramāṇa (for its effect being end of up again as an object of a pramāṇa, i.e., pratyakṣa)."
%%% THE FOLLOWING IS IRRELEVANT The big problem I have is pramāṇavyāpārānabhyupagamāt. Does this refer to a doctorinal position that does not accept the utility of pramāṇa or a consequence of the doctorinal position being discussed? It is hard to conceive anyone who does not accept the working of pramāṇa. It should be then that if we accept certain premises held by the opponent, it becomes inevitable that they also do not accept the utility of pramāṇa.
तदभ्युपगमे चानवस्थाप्रसंगः।
यत्फलावग्रहणाय प्रमाणमभ्युपगम्यते तस्याप्यन्यत्फलं तत्प्रमाणस्याप्यन्यदिति न क्वचिदवतिष्ठेत॥
\edn{tad here is pramāṇavyāpāra tasyāpīti tatpramāṇasyāpi so, pramāṇa P1 -> phala Ph1-> vyāpāra V1 is the basic. If phala is cittavṛtti grasped by pratyakṣa then pramāṇa P1 -> (cittavṛttibodha Ph1 = Ph2 <- pratyakṣa P2) -> vyāpāra V1 then pramāṇa P1 -> ((cittavṛttibodha Ph1 = Ph2 <- pratyakṣa P2) = Ph3 <- pratyakṣa P3) -> vyāpāra V1 then}%
%%% THE FOLLOWING IS IRRELEVANT % From this, it does appear that the second alternative above was the case. The opponent did not say they do not accept the utility of pramāṇa.
%%% THE FOLLOWING IS IRRELEVANT % The logic goes something like, if we accept the premise of the opponent (cittavṛtti is a) saṃvedya or b) asaṃvedya by pratyakṣa),
%%% THE FOLLOWING IS IRRELEVANT % I have the feeling this is OK by our author
%%% THE FOLLOWING IS IRRELEVANT % Oh, this is not about pramāṇaphala being the pauruṣeyaś cittavṛttibodhaḥ but about sukha, dveṣa, etc., being pratyakṣa. Our aurhot is attacking the view that does not accept those as pratyakṣa. However, I have difficulty understanding the relationship between that part (sukha, dveṣa, etc., being pratyakṣa) and the pramāṇaphala being cittavṛttibodha part.
%%% THE FOLLOWING IS IRRELEVANT % The above is probably wrong. This is not about sukhaduḥkhadveṣa, etc., but is about pramāṇa, viparyaya, smṛti, nidrā, etc.
\edn{Is Sureśvara dealing with the same type of person? ./6\_sastra/3\_phil/vedanta/naiskarmyasiddhi\_incomplete.txt:na ca sukhaduḥkhādisaṃbandho .avagatyātmanaḥ pratyakṣādipramāṇair gṛhyate yena virodhaḥ pratyakṣādipramāṇair udchāṭyate| katham| śṛṇu| EVen Dharmakīrti explicitly(?) says: rāgadveṣasukhaduḥkhādīnāṃ sarvacittacaittānām ātmasaṃvedanaṃ pratyakṣam avikalpatvāt | Cf. this from the Brahmasiddhi: tathā hi---pramāṇasya phalam arthāntaram anarthāntaraṃ vā sarvaparīkṣakaiḥ pratijñāyate prajñāyate ca| na ca tad asaṃvedyam, tadasaṃvedyatve sarvāsaṃvedyatvaprasaṅgāt, tatkṛtatvāt saṃvedyabhāvasya bhāvānām| na ca saṃvedyam, phalāntarānupalabdher anavasthāprasaṅgāc ca| tasmāt saṃvedyam, ātmaprakāśatvāt; asaṃvedyaṃ ca, viṣayavat karmabhāvābhāvāt| Maṇḍana's position is that pramāṇaphala is both saṃvedya and asaṃvedya.}%

%[25,8]
अथ न
\vrt{\dag प्रमाणव्याप\mde ारः त\ldots स्यात्\dag}{L}{%
	\rdg{\lost{5}आरः \ccs त\cce स्यात्}{\Tm},
	\rdg{प्रमाणव्यापारस्तस्य}{\Me}}% L10v7 after atha na pramāṇa
॥
\edn{The original reading in the common ancestor was "pramāṇavyāpāraḥ tasyāttadā." L actually reads: L10v7 starts vyāpāraḥ ta syāttadāpramāṇābhāvāt phalasyāsatvaṃ prasaktaṃ It is probably reasonable to assume something is missing between the ta and syāt. Since vyāpāraḥ ends with a visarga, there was an awareness that it ended an utterance. Then what would follow is tasya...āt. Then the phrase may be followed by another sentence There is probably a significant portion missing!}%

तदा प्रमाणाभावात्फलस्यासत्त्वं प्रसक्तम्। यत्र हि सदुपलम्भ\-कानि प्रमाणानि न व्यापारं गच्छन्ति, तन्नास्तीत्य्
\vrt{अभ्युपगच्छामस्तेषाम्}{L}{%
	\rdg{अभ्युपगच्छामः (तेषाम्)}{\Me}}%
॥
\edn{tad here must be phala, teṣāṃ pramāṇānām. The basic flow of events appears to be pramāṇa -> (pramāṇa)phala -> (pramāṇa)vyāpāra. The removal of teṣām from the text might be preferable but then where did it come from?}%

%[25,11]
अथागृह्यमाणमपि तदेव ज्ञा\mde नम्, प्र\mds माणं फलस्यापीति चेत्---%
% L10v8 after pramā
\edn{This does sound strange. "phalam api"? tad should be phala. "Even when not being grasped, the same [phala] is the knowledge; the phala, too, has a pramāṇa" or rather the second part should be "the pramāṇa is the evidence (pramāṇa) also for the [existence of] the phala [i.e., there is no need for the phala/jñāna to be grasped. We know that it exists because vyāpāra is being produced]" This could be an answer to the anavasthā problem in the sense there is no need for the phala to be grasped Looks like pramāṇa -> phala = jñāna -> vyāpāra}%
न, अ\mde भिव्यञ्जकधर्मव्यतिक्रमात्।
\edn{Be reminded that abhivyañjakadharma belongs to cittavṛtti according to our author}%
\vrt{न ह्यभिव्यञ्जकमगृह्यमाणम्}{\Tmpc}{%
	\rdg{न ह्यतिव्यञ्जकमगृह्यमाणम्}{\Tmac}, % Tm10v6 starts at ṇam
	\rdg{न ह्य\il{}भिव्यञ्जकमगृयमाणम्\ir{}}{L},
	\rdg{न ह्यभिव्यञ्जकं गृह्यमाणम्}{\Me}}
\edn{If bhi and ti looked similar, this is another indication that the text was from a northern script origin.}%
\vrt{अभि\-व्यञ्जयद्}{\Lpc\Me}{%
	\rdg{अभिव्य[ञ्ज\mist{1}]द्}{\Tm},
	\rdg{अभिव्य\ccs ञ्जक\cce यद्}{L}}%
दृष्टम्।
\edn{Revealer is always grasped. See above! Reminds me of the lamp being a bad example of. But isn´t this complete opposite of what our author argued for? cittavṛtti is saṃvedya/gṛhya for being abhivyañjaka? This is true if we follow the reading of Lpc (and M, A, and E)}%
न ह्य् \vrt{उल्का}{\Me}{\rdg{उत्का}{\Tm L}} गिरिशिखरेऽपि निजन्धर्म्मं अतिक्रामति॥
\edn{Perhaps the most plausible understanding of this example is that fire does what it does with or without anyone watching.}%

%[25,14]
अथाभिव्यञ्जयदेव फलं प्रत्यपि प्रमाणमिति चे\mde त्
\edn{"But the same [pramāṇa] revealing [the object] is the pramāṇa also toward the phala." "abhivyañjayat" should modify "pramāṇam" that should be the pramāṇa for phala, too -> the pramāṇa for the phala was the problem}%
सिद्धं \vrtsm{तर्हीप्सित}{\Tmpc L\Me}{\rdg{तर्हि स्मित॰}{\Tmac}}मस्माक\mds म्।
\vrt{रागादीन\mde ां}{\Me}{% L10v9 after dī
	\rdg{रगादीनां}{L}
}
वृत्तिविशेषणत्वेन ग्रहणोपपत्तेः।
\edn{Are rāga, etc., here refer to the phala? The word grahaṇa/graha is being used for the phala in this context. Be reminded that rāga and dveṣa give rise to upādāna and hāna/tyāga Need to read this a bit more carefully but looks OK in essence.}%
रक्तोऽह\mde न्द्विष्टोऽहम्
\vrt{इति च}{\Tm\Lpc}{%
	\rdg{इति चे}{\Lac},
	\rdg{इति}{\Me}}%
। यद्य्
\vrtsm{अहम्प्रत्यय}{\Me}{%
	\rdg{अहंप्रत्ययि}{\Tm L}}%
विशेषणं,
\vrt{तदापि}{\Me}{%
	\rdg{तथापि}{\Tm L}}
\vrtsm{वृत्तिभावभावि}{\eme}{%
	\rdg{वृत्तिभावाभावि}{\Tm L\Me}}%
\mds त्वेनै\-व ग्रहणाद्वृत्तिविशेषण\mde ाना\mds ं
\vrt{राग\mde ादीनां}{\Tm\Me}{%
	\rdg{\om}{L}}\edn{This is weird. Where did the edition get the reading?}
ग्रहणमुपपद्यते। \vrt{नान्यथा}{\Tm L}{% L10v10 after nyatha
	\rdg{अ(तेना)न्यथा}{\Me}}%
। रूपादि\vrtmm{वत्प्रमाना}{\Tm L}{\rdg{वन्माना}{\Me}}न्तरापेक्षत्वप्रसङ्गात्। ततश्चानवस्थेति॥
\edn{Here again, iti is used very effectively.}%

%[25,19]
\subsubsection{Final word on \emph{pramāṇaphala}}
\tstwll{PNDT}{तस्मात्\ldots उपेक्षणीयम्}%
तस्मात्प्रमाण\vrtmm{फलाभिज्ञानाभि}{L}{%
	\rdg{॰फलज्ञानाभि॰}{\Tmac},
	\rdg{॰फलानभिज्ञानामभि॰}{\Me}}%
मानमात्रमेव केवलं पण्डितम्\vrtms{मन्या\-नाम्}{\Tm\Me}{%
%	\rdg{manyānām}{\Tm\Me},
	\rdg{मन्यमानाम्}{L}}
इत्युपे\mds क्षणी\mde यमिति॥\edn{Tm has a section marker here. So, this was a major end of a discussion.}
\donote{PNDT}{Cf.\ BSSBh intro.: कथं पुनरविद्यावद्विषयाणि प्रत्यक्षादीनि प्रमाणानि शास्त्राणि चेति।
%
उच्यते---देहेन्द्रियादिष्वहंममाभिमानरहितस्य प्रमातृत्वानुपपत्तौ प्रमाणप्रवृत्त्यनुपपत्तेः। नहीन्द्रियाण्यनुपादाय प्रत्यक्षादिव्यवहारः संभवति। न चाधिष्ठानमन्तरेणेन्द्रियाणां व्यवहारः संभवति। न चानध्यस्तात्मभावेन देहेन कश्चिद्व्याप्रियते। न चैतस्मिन्सर्वस्मिन्नसति असङ्गस्यात्मनः प्रमातृत्वमुपपद्यते। न च प्रमातृत्वमन्तरेण प्रमाणप्रवृत्तिरस्ति। तस्मादविद्यावद्विषयाण्येव प्रत्यक्षादीनि प्रमाणानि शास्त्राणि च।
%
पश्वादिभिश्चाविशेषात्। यथा हि पश्वादयः शब्दादिभिः श्रोत्रादीनां संबन्धे सति, शब्दादिविज्ञाने प्रतिकूले जाते ततो निवर्तन्ते, अनुकूले च प्रवर्तन्ते; यथा दण्डोद्यतकरं पुरुषमभिमुखमुपलभ्य मां हन्तुमयमिच्छतीति पलायितुमारभन्ते, हरिततृणपूर्णपाणिमुपलभ्य तं प्रत्यभिमुखीभवन्ति; एवं पुरुषा अपि व्युत्पन्नचित्ताः क्रूरदृष्टीनाक्रोशतः खड्गोद्यतकरान्बलवत उपलभ्य ततो निवर्तन्ते, तद्विपरीतान्प्रति प्रवर्तन्ते, अतः समानः पश्वादिभिः पुरुषाणां प्रमाणप्रमेयव्यवहारः। पश्वादीनां च प्रसिद्धोऽविवेकपुरःसरः प्रत्यक्षादिव्यवहारः, तत्सामान्यदर्शनाद्व्युत्पत्तिमतामपि पुरुषाणां प्रत्यक्षादिव्यवहारस्तत्कालः समान इति निश्चीयते। (BSSBh NSP ed pp. 40–43)}
\edn{Again, the combination of tasmāt and iti So, after yeṣām was the guys who do not accept pramāṇaphala according to the 1952 edition but not according to the 1999 edition...}%
\edn{tatkāla BAUBh yathā vijñātadigvibhāgasyāpyakasmāddigviparyayavibhramaḥ. samyagjñānavato'pi cet pūrvavadviparītapratyaya utpadyate, samyagjñāne'pyavisrambhācchāstrārthavijñānādau pravṛttirasamañjasā syāt, sarvaṃ ca pramāṇamapramāṇaṃ saṃpadyeta, pramāṇāpramāṇayorviśeṣānupapatteḥ. etena samyagjñānānantarameva śarīrapātābhāvaḥ kasmādityetatparihṛtam. jñānotpatteḥ prāk ūrdhvaṃ tatkālajanmāntarasaṃcitānāṃ ca karmaṇāmapravṛttaphalānāṃ vināśaḥ siddho bhavati phalaprāptivighnaniṣedhaśrutereva; 'kṣīyante cāsya karmāṇi' 'tasya tāvadeva ciram' 'sarve pāpmānaḥ pradūyante' 'taṃ viditvā na lipyate karmaṇā pāpakena' 'etamu haivaite na tarataḥ' 'nainaṃ kṛtākṛte tapataḥ' 'etaṁ ha vāva na tapati' 'na bibheti kutaścana' ityādiśrutibhyaśca; 'jñānāgniḥ sarvakarmāṇi bhasmasātkurute' ityādismṛtibhyaśca.}%
\edn{An interesting tidbit from the Pramāṇasamuccaya 1.8 of Dignāga: tām upādāya pramāṇatvam upacaryate nirvyāpāram api sat. tad yathā phalaṃ hetvanurūpam utpadyamānaṃ heturūpaṃ gṛhṇātīty kathyate nirvyāpāram api, tadvad atrāpi. This is about pramāṇaphala. Steinkellner's edition pp. 3–4.}%

%%%%%%%%%%%%%%%%%%%%%%%%%%%%%%%%%%%%%%%%%%%%%%%%%%%%%%%%%%%%%%%%%%%%%%%%
%%%% E 25,21 -- 30,7: the Bhāṣya on anumāna and its Vivaraṇa.
%%%% Converted from the draft collation (Tm 10v9 -- 13r1; L 10v -- 13r1).
%%%% Lines marked "%% CHECK:" are decisions/readings to be verified.
%%%% Tm folio numbers are TEXTUAL (written foliation runs one behind).
%%%% Where Tm is in a lacuna (\mds ... \mde) it is not cited.
%%%%%%%%%%%%%%%%%%%%%%%%%%%%%%%%%%%%%%%%%%%%%%%%%%%%%%%%%%%%%%%%%%%%%%%%
\section{The Bhāṣya's definition of \emph{anumāna}}
%[25,21]
% Tm10v9; L10v (end)
\mll{ANUMDEF}%
\vrt{अनुमानम्पुन उच्यते \bht{ऽनुमेयस्य}}{\conj}{%
	\rdg{अनुमानमिति सुनोच्यते अ\il{}नु\ir{}मेयस्य}{\Tm},
	\rdg{अनुमानमपोच्यते अनुमेयस्य}{\Lac},
	\rdg{अनुमानमपुनोच्यते अनुमेयस्य}{\Lpc},
	\rdg{अनुमानमतोच्यते---अनुमेयस्य}{\Me}}
%% CHECK: lemma "puna ucyate" taken over from the draft as it stands (punar ucyate?).
\bht{तुल्यजातीयेष्वनुवृत्त} इति। % L10v2 after tulyaj  %% CHECK: draft says L10v2, but the preceding text was at L10v10; L11r1 follows below
\donote{ANUMDEF}{\textbf{अनुमेयस्य तुल्यजातीयेष्वनुवृत्तो भिन्नजातीयेभ्यो व्यावृत्तः सम्बन्धो यस्तद्विषया सामान्यावधारणप्रधाना वृत्तिरनुमानम्। यथा देशान्तरप्राप्तेर्गतिमच्चन्द्रतारकं चैत्रवत्, विन्ध्यश्चाप्राप्तेरगतिः।}}
अनुमेयस्य धर्मविशिष्टस्य। तद्\-धर्मवत्तया % Tm10v10 after dharmava
तुल्यजातीयेष्व् \vrt{अनुवृत्तः}{\Tm\Me}{\rdg{अनुवृत्तं}{L}} सद्भावमात्रेण,
\vrt{तद्विरुद्धधर्मकेभ्यो}{L}{%
	\rdg{तद्विरुद्धम्धर्मकेभ्यो}{\Tm\Me}}
%% CHECK: draft notes only L's "°viruddhadharma°"; Tm assumed to agree with E.
व्यावृत्तस्तेष्वसन्नित्यर्थः। \bht{सम्बन्ध} इति कः सम्बन्धः कयोर्वा? न हि सम्बन्धिनावनापेक्ष्य सम्बन्धः सम्बन्धताम् % L11r1 after sambandhināv a
%% CHECK: anāpekṣya (so the draft; anapekṣya?)
\vrt{इयान्नसौ}{\Tm\Me}{\rdg{इयान्नवसौ}{L}}
प्रसिद्धवच्चोच्यते \bht{य} इति॥ % Tm11r1 starts at vac (20210201-2204.pdf p. 41)

%[25,26]
ननु चानुमानस्य प्रकृतत्वाल्लिङ्गलिङ्गिनोः सम्बन्धः प्रसिद्ध एव नागमकत्वान्न केवलः
\vrt{सम्बन्धो}{\Tmpc L\Me}{\rdg{सम्बन्धे}{\Tmac}}
गमक उच्यते लिङ्गस्यैव % L11r2 after li
सम्बन्धो व्यावृत्त्यनुवृत्त्योः कः पुनरसौ? लिङ्गमेव। न हि लिङ्गादन्यद% Tm11r2 after anya
नुवृत्तं व्यावृत्तञ्च \vrt{तच्च}{\Tmpc L\Me}{\rdg{तञ्च}{\Tmac}} गमकन्ननु लिङ्गमपि नैव गमकं %[26,5]
लिङ्गिसम्बन्धानपेक्षं
\vrtsm{यदि चानपेक्षित}{L\Me}{%
	\rdg{\ccs यदि चानुपेक्षा\cce\ यदि चानेपेक्षित॰}{\Tm}}%
\vrtms{लिङ्गिसम्बन्धं}{\Tm\Me}{%
	\rdg{॰लिंगलिंगिसम्बन्धं}{L}} % L11r3 after liṃgaliṃ
गमकं यत्किञ्चिद्गमकं स्यात्, तस्मात्लिङ्गलिङ्गिसम्बन्धो वक्तव्यः॥

%[26,7]
नोक्तत्वाद् \bht{अनुमेयस्य तुल्यजातीयेष्वनुवृत्तः} इत्यादिनैवोक्तत्वान्न ह्यनुमेयेनासं\mds बद्धं लिङ्\mde गिना तत्सदृशजातीयेष्व् % Tm11r3 after anu; L11r4 after sadṛśajātīyeṣv a
\vrt{अनुप्रवर्त्तते}{\Tm L}{\rdg{अनुवर्तते}{\Me}}
%% CHECK: TmL anupravarttate adopted against E anuvartate.
भिन्नजातीयेभ्यो वा निवर्तते। तस्माद्\bht {य} इति च प्रसिद्धमेव संबन्धमनुवदति।

%[26,10]
\bht{तद्विषया} \vrtsm{एवमन्वय}{L\Me}{\rdg{एव समन्वय॰}{\Tm}}व्यतिरेकविशिष्ट\vrtmm{सम्बन्धि}{\Tm\Me}{\rdg{॰सम्बन्धसम्बन्धि॰}{L}}विषया % Tm11r4 after tadviṣa
%% CHECK: L °sambandhasambandhiviṣayā (sambandhasambandhiviṣayā = "having the relation of relata as object"?)
वृत्तिः। \mds\bht{सामान्यावधारणप्रधाना} लिङ्\mde\vrtmm{गि}{\Tm L}{\rdg{॰ग॰}{\Me}}सामान्यावधारण\vrtms{प्रधाना}{\Tm L}{\rdg{॰प्रधाना॰}{\Me}}। % L11r5 after liṅgi
\bht{अनुमानम्}।
%% CHECK: E prints "liṅga sāmānyāvadhāraṇapradhānānumānan"; TmL "liṅgi ... pradhānā anumānam" (pratīka without sandhi) adopted. Draft has "anumānan".

%[26,12]
\bht{\vrt{यथा}{\Me}{\rdg{तथा}{\Tm L}} देशान्तरप्राप्तेर्गतिमच्चन्द्रतारकञ्चैत्रवद्} इति। % Tm11r5 after iti
यद्यपि
\vrt{ल\mds क्षणा\mde भिधान\mds मा\mde त्रेण}{\Me}{%
	\rdg{लक्षणयाभिधानमात्रेण}{L}}
तदभिव्यक्तिस्तथाप्युदाहरणप्रत्युदाहरण\mds प्रदर्शनादतितरामि% L11r6 after atitarā
त्युदाहरण\mde प्रतिपादनन्तत्र \bht{गतिमच्च\mds न्द्रतारक\mde म्} इति चन्द्रतारकस्य गतिमत्त्वमनुमेयन्तत्समान\-% Tm11r6 after tatsamāna
जाती\mds येष्वनुवृत्ता देशान्तरप्र\mde ाप्तिर्गतिमत्सु चैत्रा\-% caitrā hard to read in Tm
दिष\mds ु, भिन्नजातीयेभ्यश्चागतिभ्यो % L11r7 after cāgatibhyo
विन्ध्यादिभ्यो व्यावृत्ते\vrtmm{त्य}{\Me}{\rdg{॰ति अ॰}{L}}व्यभिच\mde रन्ती \mds नित्यं गतिमता सम्बन्धात्।
\vrt{तस्मादित्\mde थञ्जातीयकमविनाभावेन सम्बन्धः सम्ब\mds न्ध्यन्तरं गमयति यथा वध्यघातककृ\mde दिकार्य्यकारणादि}{\Tm L}{% Tm11r7 ends after samba
	\rdg{\om}{\Me}}।
\edn{E omits this sentence, which both MSS have. Tm (partly lost) reads ...tvañjātīyakam ... sambandhaḥ samba... ; L reads sambandham for sambandhaḥ. The lacuna in Tm after avyabhicarantī covers nityaṃ gatimatā sambandhāt tasmād it.}%
%% CHECK: Tm actually reads "tvañ" for "thañ"; L "sambandham". Reconstructed lemma follows Tm's sambandhaḥ.

%[26,19]
भवतु तावदविरुद्धानान्नित्यसम्बन्धाद्देशान्तरप्राप्त्यादीनां % L11r8 after bhavatu t; Tm lacuna continues from kāraṇādi bhava
\vrt{गम्यगमकता}{\Me}{\rdg{गम्यगमक\ccs क्र्\cce ता}{L}}
वध्यघातकादीनां तु सम्\mde बन्धाभावे कथमित्युच्यते गम्यगमकभाव्भावेन।
%% CHECK: "gamyagamakabhāvbhāvena" as in the draft (gamyagamakabhāvaḥ? gamyagamakabhāvo bhāvena?).

%[26,21]
तत्रा\mds पि नित्यसम्बन्धाददोषः। न ह्यत्रापि % Tm11r8 after tatrā; L11r9 after prā
\vrt{प्राप्त\mde िलक्षणः}{\Me}{\rdg{प्राप्तिप्राप्तिलक्षणः}{L}}
सम्बन्धो\mds भिप्रेयते। सामान्यतोदृष्टानुभावप्रसङ्गान्
\edn{sāmānyatodṛṣṭānumānābhāvaprasaṅgāt?}%
न हि चन्द्रतारकादिषु देशान्तरप्राप्तेः प्र\mde ाप्तिलक्षणसम्बन्धः प्रत्यक्षेण गृह्यते। अगृह्यमा\mds णश्चेद्गमकः सर्वं सर्वस्य % Tm11r9 after agṛhyamā; L11r10 after sarvasya
गमकं स्यान्न हि किंचित्केनचिद्
\vrtsm{असम्बद्धमा}{\Me}{\rdg{असम्बमात्मा}{L}}काशादीनां सर्वत्र भावात्॥

%[26,25]
तस्मात्गम्यगमकबुद्ध्युत्पत्तिहेतुरेव सम्बन्धो अनुमानविषये \mde
तथा च बध्यघातकयोरपि अन्यतराभावद\mds र्शने % Tm11v1 after da
\vrt{ऽन्यतरभावो}{\Me}{\rdg{अन्यतराभावो}{L}} % L11r11 after nyatarābhāv
गम्यत इत्यनुमानमेव। यस्याप्युचिते देशे अभावदर्शनादन्यत्र भावो गम्यते यथा घातकदर्शनेन बध्याभावो नुमीयते चैत्रस्य जीवतो \mde % Tm11v1 ends after jīvato
गृहाभावदर्शने बहिर्भावश्चेति। % L11v1 after c

\subsection{\emph{abhāva} and \emph{arthāpatti} are \emph{anumāna}}
%[27,1]
ननु चाभावार्था\mds पत्ती एते % Tm11v3 after cābhāvārthā (so the draft; v2?)
\vrt{नानुमानलक्षणासङ्गतेर्}{L}{%
	\rdg{नानुमाने, (न) लक्षणासङ्गतेर्}{\Me}}%
गवादिवत्।
\edn{L: ete nānumāna-lakṣaṇāsaṅgateḥ, "these two are not [inferences], because the definition of inference does not fit them, as with the cow etc." E's (na) is unnecessary.}%
%% CHECK: L's reading adopted; Tm in lacuna.

%[27,2]
न य\mde थोक्तलक्षणो\mds पपत्तेः। बध्यघातकयोर्विरोधलक्षणः सम्बन्धः। यत्र यत्र घातकभावस्%
\vrt{तत्र}{\Me}{\rdg{तत्र तत्र}{L}}
बध्याभाव एवेत्यविनाभावः\mde। % L11v2 after ba
न हि बध्यस्याभावं \vrt{घातकभावो}{\Tm L}{\rdg{घातक(का)भावो}{\Me}} व्यभिचरति बध्यस्य वा\mds भावो घातकभावम्, अग्निभावमिव धूमभावः॥ % Tm11v3 after vā

%[27,6]
\mde कथं पुनरवगम्यते \mds घातकभावदर्शनादेव तद्गोचरे % L11v3 after ba
\vrtsm{बध्यस्याभाव}{\Me}{\rdg{बध्यस्याभाववे॰}{L}}बुद्धिरिति? न पुनः सदुपलम्भकप्रमाणानां घ\mde ातक इव बध्ये
\vrtsm{व्यापारा}{L\Me}{\rdg{व्यापारः}{\Tm}}\vrtms{भावादित्य्}{\Tm\Me}{\rdg{॰भादिवादित्य्}{L}}
%% CHECK: draft "vyāpāraḥ Tm" attached to vyāpārā; what follows in Tm?

%[27,8]
उच्यते। घातकभ\mds ावे व्यापृतप्रमाणस्यैव बध्याभाव\mde\vrtms{बुद्ध्युपपत्तिर्}{\Tm\Me}{\rdg{॰बुद्ध्युत्पत्तिर्}{L}}% Tm11v4 after ghātakabh
धूमभावे व्य\mds ापृतप्रमाणस्यैवाग्निभावबुद्ध्युत्पत्तिस्तथा घातकविशेषे व्यापृतप्रमाणस्य % L11v4 after vy
बध्यवि\mde शेषाभावबुद्ध्युत्पत्तिर्द्धूमविशेषव्यापृतप्रमाणस्यै\mds वाग्निविशेषबुद्धिरि\mde त्य् % Tm11v5 after pramā

%[27,12]
अभावप्रमाणवादिनस्% L11v5 after abhā
\vrt{तु सर्वत्र}{L\Me}{\rdg{तु सर्व्व्\ccs स्म\cce त्र}{\Tm}}\mds % "tu sa" hard to read in Tm
प्रमाणाभावस्याविशिष्टत्वाद्विरोधिविशेषदर्शनानुविधानं विरोध्यन्तराभावदर्शन\mde स्यानुपपन्नं लिङ्गविशेषानुविधानमिव % L11v6 after li
\vrtsm{लिङ्गिविशेष}{L\Me}{\rdg{लिङ्गिविशेषस्य}{\Tm}}द\mds र्शनस्य। \mde % Tm11v6 after liṅgiviś
\vrtsm{न हि सदुप}{\Tm\Me}{\rdg{न हि दुसदुप॰}{L}}लम्भकप्रमाणव्यापाराभावस्य
\vrt{कश्चि\mds द्विशेषो}{\Me}{\rdg{क\ccs श्चि\cce स्यचिद्विशेषो}{L}}
येन विशेषबुद्धिस्तद्विशेषमनुविधीयते।

%[27,16]
\vrt{अथ जिज्ञासिते}{\Me}{\rdg{\om}{L}}
%% CHECK: draft has "jajñāsite" (E?).
विशेष\mde े सदुपलम्भकप्रमाणा\-% L11v7 after pramāṇā
भावादिति चेत्तन्नासतः स्मृ\-% Tm11v7 after smṛ
तीच्छादिसहभावापेक्षानुपपत्तेर्न्न ह्यसत्सहायापेक्षया कार्य\mds न्निर्वर्तयद्दृष्टं शशविषाणा\mde दि दृष्टमिहेति चेन्न विप्रतिपत्तिस्थानत्वे
%% CHECK: draft "sahāyāpekṣaya"
\vrt{दृष्टान्ताभावात्}{L\Me}{\rdg{दृष्टान्तवात्}{\Tm}}% L11v8 after dṛṣṭānt
तुल्यमिति चेन्न घटादेः सतोऽपेक्षादर्शनात्॥

%[27,20]
न चापि निर्व्यापारस्यावस्तुत्वादभावस्य प्रमाणभावो युक्तः। % Tm11v8 after na c
न हि निरस्तव्यापारकं किंच\mds ित्कार्या\mde\vrtms{रम्भकं}{\Tm\Me}{\rdg{॰रम्भन्}{L}}दृश्यते लोके। % L11v9 after lok
तस्माद्व्यापाररूपप्रमाणजनिता प्रमेयाभावबुद्धिः
\vrt{प्रमेयविषयत्वात्}{\Tmpc\Me}{%
	\rdg{प्रमेयविषयत्वा\il{}त्\ir{}}{\Tm},
	\rdg{प्रमेय\ccs वि\cce विष\ccs य\cce विषयत्वात्}{L}}%
घटादिबुद्धिवत्। तथा च वस्तुत्वात्प्रवृत्तिनिवृत्तिहेतुत्वाद्विपक्षः शशविषाणादिस्% Tm11v9 after va
\vrt{तथा अभावप्रमाणभावजनिता}{\Tm L}{% L11v10 after pramā
	\rdg{तथाभावप्रमाणभावजनिता}{\Me}}
%% CHECK: TmL "tathā abhāva°" (hiatus) adopted.
न भवति, यथाभिहित\vrtmm{हेतु}{\Tm\Me}{\rdg{॰हेत॰}{L}}दृष्टान्तदर्शनेन।

%[27,25]
घातकदर्शनं बध्यभावबुद्धेर्लिङ्गं
%% CHECK: badhyabhāvabuddher (so the draft; badhyābhāva°?)
\vrt{तद्भावभावित्वाद्}{\Tmpc L}{%
	\rdg{तद्भावभाविस्वाद्}{\Tmac},
	\rdg{तद्भावित्वाद्}{\Me}}
%% CHECK: TmL "tat bhāvabhāvitvād" adopted (E tad bhāvitvād).
अग्न्यादिबुद्धिनिमित्तधूमादिबुद्धिवत्% Tm12r1 starts here (bottom strip, 20210201-2204.pdf pp. 41/43; left end abraded)
\vrt{संशयनिवर्तकत्वाच्च}{\Tmpc\Me}{%
	\rdg{संशयनिवर्त्तकात्वाच्च}{\Tmac},
	\rdg{संशयनिवर्त्तकाच्च}{L}}
\vrtsm{बध्या}{L\Me}{\rdg{व्यव्या॰}{\Tm}}भावबुद्धिर्वा लिङ्गजनिता निश्चयरूपत्वादग्निबुद्धिवदिति। % L12r1 after bhāva
%% CHECK: Tm reads "vya(?)" for ba and "vyā" for dhyā.

%[28,1]
इह घटो नास्तीति किल नानुमानं लिङ्गलिङ्गिसम्बन्धाग्रहणान्न हि यथाग्न्यादिबुद्धेर्धूमादिलिङ्गान्वेषणं तथेह नास्ति घट % Tm12r2 after ghaṭa
इत्यत्र लिङ्गान्वेषणमस्ति प्रागेव लिङ्गान्वेषणात्बुद्ध्युत्पत्तेर्। % L12r2 after a

%[28,4]
नैष दोषः। इहप्रत्ययाकारस्यैव लिङ्गत्वान्न हि
\vrt{प्रथममेव}{\Tmpc\Me}{%
	\rdg{प्रथनिमव}{\Tmac},
	\rdg{प्रथमेव}{L}}
लिङ्गबुद्ध्युप\vrtms{जनने}{L\Me}{\rdg{॰ज केन}{\Tm}}
सति लिङ्गान्वेषणं भवतीहेत्याकारोपरञ्जिता बुद्धिर्ल्लिङ्गम्
\vrt{उपरञ्जक}{\Tmpc\Me}{\rdg{उपरञ्ज}{\Tmac L}}\-% L12r3 after vyati; Tm12r3 after rikta
व्यतिरिक्तघटाद्यभावबुद्धेर्घातकबुद्धिवदेव विरोधलक्षण एव चात्रापि सम्बन्धः।

%[28,8]
य\mds त्र पुनरविरुद्धान्युपर\mde ञ्जकानि तत्रास्तित्वबुद्धेः।
\vrt{एतेनैवार्थापत्तेर्}{\Tmpc L\Me}{\rdg{एतेनैवार्थावत्तेर्}{\Tmac}}
अनुमानत्वं व्याख्यातम्
\vrt{अविनाभावेन}{\Tmpc L\Me}{\rdg{अपिनाभावेन}{\Tmac}} % L12r4 after avin
सम्बन्धेन तत्रापि जीवतो हि % Tm12r4 after jīvato
\vrt{देवदत्तस्य}{\Tmpc\Me}{%
	\rdg{देवात्तास्य}{\Tmac},
	\rdg{देवात्तस्य}{L}}
गृहाभावे बहिर्भावः, यदा गृहे भावो न तदा बहिर्भाव \mds इत्य्
\vrtsm{अन्यतरेणानु}{\Me}{\rdg{अन्यतरेणान्यतरेणानु॰}{L}}मानम्। सा\mde क्षात्दृश्यमानबहिर्भावाभ्यन्तराभावत्वन्दृष्टान्तः। % L12r5 after abhyanta
%% CHECK: draft has \mde both after sā and after dṛśya; the first is kept.
वृत्त्युपरञ्जकानुपरञ्जकत्वेन ह्यभावग्रहणस्याभिहितो हेतुर्। % Tm12r5 after abhihit

%[28,13]
अग्न्य\mds ादावपि धनादिश\mde क्तिग्रहणमनुमानमेव।
%% CHECK: dhanādi° (so the draft; dahanādi°?)
कार्यन्निर्व्व\mds र्तकयोरविनाभावसम्बन्धस्यात्मनि
\vrt{दृष्टत्वाद्}{\Me}{\rdg{दृष्ट\ccs मि\cce त्वाद्}{L}}% L12r6 after dṛṣṭa
भुज्या\mde दौ। दृश्यते हि
\vrt{प्रत्यगात्मनि}{\Me}{\rdg{प्रत्यात्मनि}{L}}
%% CHECK: draft has \mde both after bhujyā and after śakt; Tm's state for "dau ... śakt" is unclear, so Tm is not cited here.
शक्तो हमिदं कार्यं कर्त्तुं अशक्तो हमित्यात्मवदात्मविशेषण\mds स्य शक्तेर्ग्रहणं शक्तिनिर्व\mde र्त्त्यमेवन्दहनादिकार्य\mds मपीत्यनुमानसिद्धिः। % Tm12r6 after ātmaviś

%[28,17]
ननु चाभावस्य प्रत्यक्षतामुपगायन्ति केचिद्। यथा किलादर्शकेन प्रदीपादि\mde ना सति गृह्यमाणे वस्तुनि न तदवयविनि यदुप\mds लभ्यते तन्नास्तीत्येवं सतः % L12r7 after nanu c (draft: r6); Tm12r7 after upa
\vrt{प्रकाशकं}{L}{\rdg{प्रकर्शकं}{\Me}}
प्र\mde माणमसदपि प्रका\mds शयति। % L12r8 after śaya (draft: r7)
तथा चोदाहरन्त्यनङ्कितं वास आनयेत्युक्तः प्रविश्य भाण्डागारमंकितानङ्कितवासः सन्नि\mde धावंकितवासो ग्रहणेनैव प्रत्यक्षेण\mds ाङ्काभावविशिष्टं वास उपलभ्यनयति। तस्मात्\mde% Tm12r8 after pratyakṣeṇ
%% CHECK: \mds after pratyakṣeṇ supplied (draft has \mde twice in a row: after sanni and after tasmāt). upalabhyanayati (upalabhyānayati?).
प्रत्यक्षमेवाभा\mds वो लक्षणान्तरानुपपत्तेरिति॥ % L12r9 after lakṣaṇ

%[28,23]
तन्न।
\vrtsm{बुद्ध्याकारेणा}{\Me}{\rdg{बुद्ध्याकोरेणा॰}{L}}र्थावधारणानृतत्वप्रसङ्गाद्यदाकारा हि बु\mde द्धिस्तस्य प्रत्यक्षावध्रियमाणात्मताभ्युपगम्यते % Tm12r9 after ya
\vrt{य\mds था}{\Lac\Me}{\rdg{यदा}{\Lpc}}
नीलं वस्तु पीतमित्यत्रापि विशेषणविशेष्यसम्बन्धप्रत्ययो मिथ्या प्राप्नोति। % L12r10 after prāpnoti
तद्यथेह घटो नास्तीतीहप्रदेशस्य घटाभावेनासति सम्बन्धे सम्बन्धप्रत्ययः \mde
न हि \vrtsm{घटाद्याभाव}{\Tm\Me}{\rdg{घटाभाव॰}{L}}स्येहप्रदेशस्य च प्राप्ति\mds लक्षणः समवायाभिधानो वा सम्बन्धस्तैरभ्युपेयते यथेह देवदत्त इति। % Tm12v1 after prāpti
%% CHECK: draft reads "prāpti%Tm12v1 \mds prāptilakṣaṇaḥ" -- dittography in E or in the draft?
तस्माद् % L12v1 after abhā
\vrt{अभावप्रत्ययवन्}{\eme}{%
	\rdg{अभावे प्रत्यया}{L},
	\rdg{अभाव(वे)प्रत्ययवा(या)न्}{\Me}}
\edn{abhāvapratyayavat, "like the cognition of absence": E's parentheses spoil the point.}%
%% CHECK: emendation; the draft's note for L is "abhāve pratyayā L".
नीलं वस्त्विति विशेषणविशेष्यसम्बन्धप्रत्ययो अपि सम्बन्धप्रत्ययत्वान्मि\mde थ्या प्रसज्येत।

%[29,4]
यथा \vrt{चास्त्य्}{\Tmpc L\Me}{\rdg{चावस्त्य्}{\Tmac}} अपीहदेशघटाभाव\mds % Tm12v2 after ghaṭābhāva
\vrt{सम्बन्धे सम्बन्धप्रतीतिस्}{\Me}{\rdg{सम्बसबन्धप्रतीतिस्}{L}}%
%% CHECK: draft E text "sabandha" corrected to sambandha.
तथैवानीले पि
\vrt{वस्तुन्यसत्येव}{\Me}{\rdg{वस्तुन्न्येवसत्येव}{L}}
\vrt{नीलत्वसम्बन्धे}{\Me}{\rdg{नीलत्वे सम्बन्धे}{L}} % L12v2 after nīla
नीलसम्बन्धप्रत्ययः प्राप्नोति, तथा सम्बन्धिष्वप्यथास्त्येव सम्बन्धः द्विष्ठता विरुद्ध्येत सम्बन्धस्य। अ\mde भावस्यावस्तुत्वाभ्युपगमाद्।

%[29,7]
अथाप्येकशतं षष्ठ्यर्था इ\mds ति यावत्षष्ठ्यर्थाः सम्बन्ध इष्येत, त\mde दा द्विविध एव सम्\mds बन्धः समवायाख्यः सम्योगाख्यश्चेति न युज्यते॥ % Tm12v3 after i; L12v3 after iṣy

%[29,9]
\vrtsm{तथैक}{L}{\rdg{अथै(तथै)क॰}{\Me}}शतस्यापि षष्ठ्यर्थानां % L: ta cancelled before tathaika
\vrt{उभयान्तःपातित्वम्}{\Me}{\rdg{उभयान्तःप्र्\mist{2}तित्वम्}{L}}
उच्येत \vrt{नाभावे}{\Lac\Me}{\rdg{न भावे}{\Lpc}}
\vrt{ऽन्यतरस्याप्यसम्भवात्। किञ्चान्यत्। इह}{\eme}{%
	\rdg{न्यतरस्याभावानित्वन्यदिह}{L},
	\rdg{ऽन्यतरस्याप्यसम्भवात्॥ किञ्चा(न्त्व)न्यत्। इह}{\Me}}
%% CHECK: L's reading as in the draft; E prints kiñcā(ntva)nyat.

%[29,11]
नभसि कुसुमं नास्तीति सम्बन्धिनो % L12v4 after na (draft: v3)
\vrtsm{नभस\mde ो अप्यप्रत्यक्ष}{\Me}{%
	\rdg{नभसो प्यक्ष॰}{\Tm},
	\rdg{नभसो प्रत्यक्ष॰}{L}}त्वात्गगन\mds कुसुमाभावस्याप्रत्यक्षत्वम्प्राप्नोत्यथाप्रत्यक्षत्वमेवेति
\vrt{चेदापतितं}{\Me}{%
	\rdg{चेदावितः}{\Lac},
	\rdg{चेदावतः}{\Lpc}}
गगनकुसुमाभाव\mde वदप्रत्यक्षत्वमिह भूप्रदेशे घटाभावस्यापि। विशे\mds षो वक्तव्यो। % L12v5 after va (draft: v4); Tm12v5 after ghaṭābhāvasyā (draft: 13v5)

%[29,14]
अथ व्यो\mde म्नोऽप्रत्यक्षत्वादिति तन्न। पुष्पस्य\mds पततो नभस्यस्तित्वदर्शनाद्व्योमाप्रत्यक्षत्वमहेतुः। यथैव निपतत्कुसुमसद्भा\mde वे % L12v6 after yathai
\vrt{व्योमाप्रत्यक्षत्वमहेतुस्}{L\Me}{\rdg{व्योमाप्रत्यक्षत्वहेतुस्}{\Tm}}%
तथा सद्भाव\vrtmm{विरुद्ध}{L\Me}{\rdg{॰विरुद्द्\ccs\mist{2}\cce द्ध॰}{\Tm}}ग्रहणे प्यहेतुः स्यादथाकाशाप्रत्यक्षत्वादेवेति चेदसम्ब\mds न्धमपि यत्किंचिद्व्यवहिताप्रत्यक्षत्वादि सो अपि हे\mde तुः स्यादथ % Tm12v6 after gra; L12v7 after s
\vrt{नभःकुसुमाभावः}{\Tmpc L\Me}{\rdg{नभःकुसुमावः}{\Tmac}}
प्र\mds त्यक्ष इ\mde ति चेदिन्द्रियार्थसन्निकर्षानर्थक्यं सर्वत्र
\vrt{प्रसज्येत}{\Tm L}{\rdg{प्रसज्यते}{\Me}} % Tm12v7 after prasajya
%% CHECK: TmL prasajyeta adopted.
तथा च सत्यसन्निकृष्टोपलब्धेः सर्वेषां सर्वज्ञत्वं प्रसक्तम्।

%[29,20]
अ\mds थानुमेयन्तत्रेति चेन्नभः कु\mde सुमाभाववद् % L12v8 after ce
\vrt{इह घटाभावस्याप्यभावाविशेषाद्}{\Tm}{%
	\rdg{इह घटाभावस्याप्य\ccs भा\cce वाविशेषाद्}{L},
	\rdg{इह घटाभावस्याविशेषाद्}{\Me}}
%% CHECK: the draft's second reading has no siglum; taken as L.
\vrt{आनुमानिकीयमेव}{L}{%
	\rdg{आनुमानकीयमेव}{\Tmpc},
	\rdg{आनुमानक्रीयमेव}{\Tmac},
	\rdg{अनुमानिकत्व(कीय)मेव}{\Me}}
\edn{ānumānikam eva?}%
सिद्धम्। एवञ्च सत्यभावप्रमाणवादिनो अपि व्यवहितविप्रकृष्टोदकादावनलाद्यभावस्य्\-% Tm12v8 after vādin
\vrt{आनुमेयत्वादिह}{\Tmpc L\Me}{\rdg{आनुमेयत्वाह}{\Tmac}}
\vrtsm{घटा}{\Tmpc L\Me}{\rdg{घा॰}{\Tmac}}भावो अपि तत्सा\mds मान्या\mde दनुमानग्राह्य एवेति % L12v9 after anumā
\vrt{निश्चेतुं युक्तम्}{L\Me}{\rdg{निश्चेतुं \il{}\il{}युक्तं\ir{}ऽ}{\Tm}}।
%% CHECK: the draft's ``yuktaṃ'' for Tm reproduced as is (marked in the MS?).

%[29,24]
\bht{\vrt{विन्ध्यश्चाप्राप्तेरगतिरिति}{\Tmpc\Me}{%
	\rdg{विन्ध्यश्चाप्राप्तेरगिरीति}{\Tmac},
	\rdg{विन्ध्यश्चाप्राक्तेरगितिरिति}{L}}}
च प्रयोगान्तरम्। अगतिर्विन्ध्य इति प्रतिज्ञानम्। अप्राप्तेरिति हेतुश्% Tm12v9 after he
\vrt{चैत्रवद्}{\Tm\Me}{\rdg{चेत्रवद्}{L}}
इति काकलोचनवदुभयत्र सम्बध्यते। गतिमच्चन्द्रतारकमित्यनेन साधर्म्येण विन्ध्यश्चाप्राप्तेरित्यनेन % L13r1 after vindhyaś cā
\vrt{वैधर्म्येण}{\Tm\Me}{\rdg{वेधर्म्म्येण}{L}}
सम्बन्ध। एवञ्च सत्य्
\vrt{अन्वयव्यतिरेकदर्शनकृतोऽदोषो}{\Tmpc L}{%
	\rdg{अन्वय\ccs न्ध्य\cce व्यतिरेकेणदर्शनकृतो दोषो}{\Tmac},
	\rdg{अन्वयव्यतिरेका(क)दर्शनकृतो न दोषो}{\Me}}
\edn{TmL read darśanakṛto doṣo without na, which is to be read darśanakṛto 'doṣaḥ (cf. adoṣaḥ at the end of the next sentence). E's (ka) and na are unnecessary.}%
%% CHECK: Tmac "tirekeṇa" and Tmpc "tireka" are both marked uncertain in the draft.
%[30,7]
ऽथ वैतस्मिन्प्रयोगे व्यतिरेकविपर्ययदर्शनेन दोषाभाव इत्युपन्यस्तं लक्षणस्य च % Tm13r1 after tha v
\vrt{लक्ष्यविशे\mds षेणोह्यत्\mde वाद्}{\Me}{\rdg{लक्ष्यविशेषणोह्यत्वाद्}{L}}
अदोषः॥\edn{Tm has a section marker here.}%

\end{document}

[30,9] आप्तः परानुजिघृक्षाप्रयुक्तो% °kṣo Tmac
दोषरहितो दृष्टानुमितार्थः %Tm13/14r2
तेन दृष्टो वा अनुमितो वार्थः% Tm13/14r2 %tena (dṛṣṭo vānumito vārttha) cancelled in Tm; dittography
परत्र श्रोतृविशेषे स्वबोधसङ्क्रान्तये स्वाधिगमसंक्रमणाय शब्देन %\mds
शब्द इति वाक्यं त%\mde
तो विशिष्टार्थप्रतिपत्तेस्तस्मात्तदुपदिष्टाद्वाक्याद्वाक्यार्थविषया या चित्तवृत्तिः सानुमानवदेव सामान्यावधा%Tm13/14r3
रणप्रधाना

[30,14] इत्थमाप्त्यादिधर्मग्रहणपूर्वका% °ko Tmac
श्रोतुरागमः न वक्तुर्दृष्टानुमि%\mds
तार्थो हि वक्ता तस्य लैङ्%\mde
गिकमैन्द्रियकं वा तज्ज्ञानम्श्रोतुरेव तु श्रवणपूर्वकन्निश्चितं%niścita Tm
विज्ञानं लैंगिकैन्द्रियकव्यतिरिक्तं स %Tm13/14r4
आगमः न प्रत्यक्षमिन्द्रियाविषयत्वाद्यथाभिहितलिङ्गलिङ्गिसम्बन्धनि%\mds
रपेक्षत्वाच्च नानुमानम्। न हि स्वर्गा%\mde
पूर्व%\mds
देव%mde
तादिशब्दार्था%°śabdo rtthārthā° Tmac
न्वितं ज्ञानमासादितलिङ्गलिङ्गिसम्बन्धपु%°sambandhasu° Tmac
रस्सरमुपजायते। फलं तु सर्वत्र प%Tm13/14r5
उरुषेयो%\mds
वृत्तिबोध एव॥

[30,20] यस्यागमस्य %\mde
वक्तागमाख्यवृत्त्युत्पत्ति% ``. . . . ''
निमित्ते %\mde
श%\mds
ब्दोच्चारणातागमस्य प्रयोगोपन्यासेन प्रतिपा%\mde
दितं फलन्तु पूर्वव%\mds
देव यदाप्येक %\mde
एव प्रयोगस्तदापि साध्यसाधनाभाव%°sādhanabhāva° Tm
मात्रस्य विव%Tm13/14r6
क्षितत्वा%\mds
न्न विपरीतव्यतिरेकदर्शनादित्यु%\mde
च्यते

[30,23] अश्रद्धेयार्थ%\mds
वक्त्रभिहितत्वेनागम एव विशिष्यते। कथमसावश्रद्धेयार्थो वक्तेत्यत आह न दृष्टो%\mde
ऽनुमितो वार्थस्तेन वक्ता शिष्टलेकप्रसिद्धान्यथा%Tm13/14r7
 प्रतिपादनात्बुद्धार्हतादिः। न हि %\mde
रागादिदोषानु%\mds
पहतान्तःकरणः प्रसिद्धमन्यथा निर्वर्णयति। स आगमः प्लवत आभासी-भवति॥

[30,27] यद्यपि प्रत्यक्षानुमान%\mde
विषयाभासदर्शनं स्वशब्देन न कृतन्तथा अप्यागमाभासो%Tm13/14r8 \mds
पदर्शनेनैव पूर्वत्रापि यथाभिहितलक्ष%\mde
णवि%\mds
परीतत्वेन प्रत्यक्षानुमानाभासादि द्रष्टव्यमर्थादुक्तम्॥

[31,1] मूलवक्तरि त्वाद्यवक्तरीश्वरे तन्न्निमित्तो य आगमः स वि%\mde
प्लवनिमित्ताभावान्निर्विप्लवः स्यात्। यथा-अर्थ(य)वक्तृवत्॥

[31,3] उपमानं (मानं) % syā``dva''ktrupamānaṃ
शब्दपू%Tm13/14r9 \mds
र्वकं तदागमान्तःपातित्वान्न
प्रम्णान्तरम्। यथा गौरिव गवय इति वनेचरभिहतादेव वाक्यादवधृतसादृश्य पुनर्वने गवयदर्शने अपि
पूर्वश्रुतवाक्योपस्थापि%\mde
तसादृश्यस्मरणमात्रमेव नाधिकम्। सादृश्यं तु वा%Tm13/14v1 \mds
क्यादेव पूर्वप्रतिपन्नम्॥

[31,7] नापि संज्ञासंज्ञिसम्बन्धप्रतिपत्तित एव मानं तस्या अपि तद्वाक्यादेवावगतत्वात्। न हि गौरित्येतावतानवगतगवयनामकस्य वनगतगवयदर्शिनो अपि ग%\mde
वयोऽयमिति संज्ञासंज्ञिसम्बन्धावबो%Tm13/14v2 \mds
धो भवति॥

[31,10] तस्माद्गौरिव गवय इति गवयशब्दो अपि तद्वाक्येऽपेक्षितव्यः। तथा च सति पूर्वश्रुतादेव सर्वं प्रतीतम्। अथापि स्यादयं(दस्य)गवय इति न पूर्वबुद्धिविषय%\mde
मिति तदपि तत्रभवतामपि परिचोदनपरिहारमपश्यता%Tm13/14v3 \mds
मपि
पीडितकैवले[कौले]यकवृत्तिकोणोऽ%\mde
वस्थानमिव%\mds
लक्ष्यते॥

[31,14] तथा हि देशकालसन्निकर्षविप्रकर्षादिभेदेननेकोपमानभेदः प्रसज्येत। तथास्त्विति चे%\mde
त्तत्पूर्वधिया एव गतमन्तर्निहितसमस्तविशेषमेव%Tm13/14v4 \mds
हि वा नेयौ पूर्वमुपपादितम्। न हि गौ%\mde
रित्युक्ते देशका%\mds
लादिविशेषाभावे न गोशब्दादर्थप्रतिपत्तिः॥

[31,17] किं चान्यत्न हि संज्ञासंज्ञिसम्बन्धप्रति%\mde
पत्तिरुपमानं येन ह्यपमीयते तदुपमानं सा%Tm13/14v5
दृश्यज्ञा%\mds
नं हि तत्। न हि नामिनामसम्बन्धप्रतिपत्त्या किञ्चिदु%\mds
पमीयत उत्तरकालत्वात्। अवधृतसादृश्यस्य हि पश्चादा(स्य)यं गवय इत्येवं सो अपि निपत%\mde
ति

[31,21] एवन्तर्हि स्यादशब्दपूर्वकं प्रमाणान्तरं यथा गव%Tv13/14v6
यदर्शनादे%\mds
तत्सदृशो गौरिति ग्रा%\mde
मीणस्य बुद्धिस्तदपि न
प्र%\mds
मेयाभावात्। सति हि प्रमेये सदुपल्म्भकत्वा%\mde
भिमतस्य प्रमाणस्य व्यापारः न हि गोगवययोर्व्यावृत्तमनुवृत्तं च%vyāvṛttānu<pa>vṛttañ ca tm
वस्तुतः सादृश्यन्नामै%Tm13/14v7
कमस्ति॥

[31,25] ननु च सादृश्यबुद्धिरस्त्येव नासौ निरालम्बना भवितुमर्हति%\mds
सत्यमेवं तत्रैव वि%\mde
चारणा किमसौ वनयूथादिप्रत्ययवदथ घटादिबुद्धिवदिति किञ्चातः यदि वनप्रत्ययवद्भवेत्तदा%pratyayabhavet tathā
य%Tm13/14v8
थैव विशिष्टसन्निवेशानुगता महीरुहा एवात्मव्यतिरेकेणाविद्यमानवनप्रत्ययहेतवस्तथैव गोगवयावयवा एव द्वयदर्शनस्मृत्यपेक्षया स्वात्मव्यतिरेकेणासतः सादृश्यस्य
प्रतीतौ कारणमिति प्रा%Tm13/14v9
प्तमवस्तुत्वम्

[32,1] अथ पुनर्घटादिप्रत्ययवद्वस्तु%°pratyayavastu° Tmac
त्वं% tva Tmac
प्रमाणत्वञ्च तदा स्यादिति

[32,2] तत्र कीदृशं वस्त्विति विचार्यते 32,3  किं गोगवययोरनुवृत्तमेकं सादृश्यं गोत्वगवयत्वाभ्यामन्यत्% anya Tmac
? ते एव वान्योन्यविशिष्टे %\mds
यदि वा %Tm15r1
ताभ्यां भिन्ने गोगवयाधारे द्वे सादृश्ये परस्परानुपेक्डे, सापेक्षे वा ? 32,6  किं व्यक्तिरेव सान्यतरजातिविशिष्टा वा किं जातिः साचान्यतरव्यक्तिविशिष्टा व्यक्तिद्वयविशिष्टा वा सम्बन्धो वैतेषां यथा-विकल्पयोगमित्येकैकान्योन्यविशिष्टविशेषणत्वादिना भूयांसोऽन्ये अपि विकल्पा उद्भावयितव्याः। तत्र यद्येकं सादृश्यमुभयगतं तदा नोपमेयत्वमेकत्वात्यथा गौरिव गौरिति॥

[32,10] अथ तेनान्यतरव्यक्तिरुपमीयते नैतदपि स्वात्मगतत्वात्गौरिव गोत्वेन। अथ तद्वत्या व्यक्त्या व्यक्त्यन्तरमुपमीयते नैतदपि न हि गोत्व्वं खण्डमुण्डोपमेयत्वकारणम्। अथ धर्मान्तरापेक्षया सादृश्यं तद्बुद्धेः कारणमिति नष्टं तर्हि सादृश्यं स एवास्तु धर्मः तद्बुद्धेः किं ते सादृश्यहतकेन कल्पितेन॥

[32,14] किं च स चाप्यपेक्ष्यमाणो धर्मः किमेकः सादृश्यवदुभयत्राथ भिन्न इत्येवमादिविकल्पाः स्युरित्यसम्भवः पूर्ववदेव। अथ भिन्नं सादृश्यमुच्यते तच्च न भेदादेव न हि भवति गौरिवाश्व इति। यथैवैकैकत्वादघटमानता तथा बहुभिरपेक्ष्यमाणैः। एतेन जातिव्यक्तिसम्बन्धादयः पूर्वविकल्पिताः
प्रत्युक्ताः॥

[32,19] अथ बुद्ध्युत्पत्त्यन्यथानुपपत्त्या सादृश्यकल्पना सत्येवं वनविवरपूर्वदक्षिणोत्तरादिबुद्धियोगान्यथानुभवात् [न्यथानुपपत्त्या] वस्तुत्वं किं न कल्प्यते। स्यादेषा तव बुद्धिर्वृक्षादय एव तत्र तद्बुद्धिजन्महेतव इति। यद्येवं मिथ्याप्रत्ययो गवयेन सदृशी गौरिति वनबुद्धिवत्। अथ ये व्यावृत्ताः क इहोपमार्थः गौरिवाश्व इति। अथोभये अपि तेऽवयवा व्यावृत्तनुवृत्तात्मानस्तद्बुध्दिहेतवः तथापि वनप्रत्ययवदेव
प्राप्तम्॥

[32,25] तस्मात्परिक्ष्यमाणं नातिविमर्दनक्षममिति त एवावयवा उभयत्रोपलभ्यमाना भूयांसः सादृश्यबुद्धेः कारणं वनबुद्धेरिव
पादपाः [इति] नोपमानं नाम प्र्माणान्तरमिति सिद्धम्॥

[32,27] नापि प्रेमेयद्वित्वात्प्रमाणद्वित्वकल्पनम्। कुतः प्रमाणवशादर्थप्रतीतेः। न ह्यर्थवशात्प्रमाणकल्पना। यदि चार्थवशात्%
प्रमाणकल्पना स्यातधुना 33,4       तन्न शक्यं जानानामसर्वज्ञत्वात्चैत्यवन्दनात्स्वर्गो भवतीत्यत्र चैत्यवन्दनस्य स्वर्गहेतुत्वबुद्धेः प्रत्यक्षनुमानाविषयत्वान्मिथ्यात्वं स्यात्। तदर्थं वा प्रमाणान्तरकल्पनम्॥

[33,7] अथ सा बुद्धिरन्मानफलमिति चेत्प्रमाणादन्यत्र
फलप्रसङ्गः एवं चनिर्मोक्षः(क)रणत्रयादिवाक्यार्थविष्यबुद्धिरपीति॥

[33,9] नापि चात्र प्रमाणलक्षणनिरूपणतात्पर्यं
प्रसिद्धप्रमाणादिवृत्तिनिरोधाभ्युपायाभिप्रदर्शनपरत्वात्। अत एव न सूत्रकारः सूत्रयाञ्चकार। भास्यकारेण तु किञ्चित्%
प्रसिद्धानुवादमात्रमवलम्बितमिति तदनुसरेणेन किमपि जल्पितमिति॥6-7॥
